% Unobservable effects: fractional permissions as a discipline for in-place
% update in Lean 4.
%
% Build, from this directory. Highlighting is minted's, so the build needs
% shell escape and Pygments on PATH:
%   python3 gen-laws.py     # the law frames, after a change to the library
%   python3 gen-frames.py   # the frames themselves: the quoted ranges, comment-free
%   latexmk -pdf -shell-escape unobservable-effects.tex
% Check that every frame still quotes the library:
%   python3 check-listings.py
%
% pdflatex: the toolchain's lualatex cannot load its font loader, so the Unicode
% in the quoted Lean is mapped to font commands by newunicodechar below. A
% character in a quoted range with no mapping fails the build, which is the
% signal to add one.
\documentclass[11pt,a4paper]{article}

% pdflatex: the toolchain's lualatex cannot load its font loader, so the
% Unicode in the quoted Lean is mapped to font commands below, both in the
% body (newunicodechar) and inside lstlisting (literate).
\usepackage[T1]{fontenc}
\usepackage[utf8]{inputenc}
\usepackage{lmodern}
\usepackage{amsmath,amssymb,amsthm,mathtools}
\usepackage{mathpartir}
\usepackage{newunicodechar}
\newunicodechar{×}{\ensuremath{\times}}
\newunicodechar{α}{\ensuremath{\alpha}}
\newunicodechar{β}{\ensuremath{\beta}}
\newunicodechar{σ}{\ensuremath{\sigma}}
\newunicodechar{τ}{\ensuremath{\tau}}
\newunicodechar{ℓ}{\ensuremath{\ell}}
\newunicodechar{→}{\ensuremath{\to}}
\newunicodechar{∀}{\ensuremath{\forall}}
\newunicodechar{∃}{\ensuremath{\exists}}
\newunicodechar{∈}{\ensuremath{\in}}
\newunicodechar{∉}{\ensuremath{\notin}}
\newunicodechar{∧}{\ensuremath{\wedge}}
\newunicodechar{≠}{\ensuremath{\neq}}
\newunicodechar{≤}{\ensuremath{\leq}}
\newunicodechar{⟨}{\ensuremath{\langle}}
\newunicodechar{⟩}{\ensuremath{\rangle}}
\newunicodechar{…}{\ensuremath{\ldots}}
\newunicodechar{¬}{\ensuremath{\neg}}
\newunicodechar{γ}{\ensuremath{\gamma}}
\newunicodechar{∨}{\ensuremath{\lor}}
\newunicodechar{↔}{\ensuremath{\leftrightarrow}}
\newunicodechar{≃}{\ensuremath{\simeq}}
\newunicodechar{₀}{\textsubscript{0}}
\newunicodechar{₁}{\textsubscript{1}}
\newunicodechar{₂}{\textsubscript{2}}
\newunicodechar{◾}{\ensuremath{\blacksquare}}
\usepackage{booktabs}
\usepackage{enumitem}
\usepackage[margin=1in]{geometry}
\usepackage{microtype}
\usepackage{xcolor}
\usepackage{minted}
% gruvbox-light: the most hues of the styles that stay legible on white paper.
% The dark styles (one-dark, monokai, dracula, lovelace) paint identifiers in
% near-white, which vanishes here; a dark editor look would need a dark
% background, which \setminted{bgcolor=...} could give if it were wanted.
\usemintedstyle{gruvbox-light}
\setminted{fontsize=\footnotesize,frame=single,rulecolor=\color{black!30},
           framesep=4pt,breaklines=true,baselinestretch=1.0}
% A quoted frame is written \codeframe[linerange=lo-hi]{../path}: a line range of
% a library file. The frame itself is generated -- the range with its comments
% removed, since a frame is code and not documentation -- into frames/<key>.lean
% by gen-frames.py, and frames/lookup.tex maps the range to that file. A call
% site therefore names the source it quotes and not the frame, and a range with
% no generated frame is a build error rather than a silently missing box.
%
% A frame is one unbreakable box, so a frame that does not fit the space left on
% a page moves to the next one instead of being split across two.
% Generated by gen-frames.py; do not edit by hand.
% Maps a \codeframe call's (path, range) to its generated frame, so a
% call site names the source and not the frame. The frames themselves are
% frames/<key>.lean, comment-free copies of the ranges, checked by
% check-listings.py against the library.

\expandafter\def\csname cf@../Scidonia/FracPerm/Ledger.lean@25-25\endcsname{9750611a7baa}
\expandafter\def\csname cf@../Scidonia/FracPerm/Ledger.lean@28-28\endcsname{5f928bda90e1}
\expandafter\def\csname cf@../Scidonia/FracPerm/Ledger.lean@34-35\endcsname{1913b63860bf}
\expandafter\def\csname cf@../Scidonia/FracPerm/StateVar/Value.lean@84-84\endcsname{dbe771f805da}
\expandafter\def\csname cf@../Scidonia/FracPerm/StateVar/Value.lean@96-97\endcsname{c011a33daef6}
\expandafter\def\csname cf@../Scidonia/FracPerm/StateVar/Value.lean@109-109\endcsname{c6b92c041b00}
\expandafter\def\csname cf@../Scidonia/FracPerm/StateVar/Value.lean@113-113\endcsname{9c6db65f9b82}
\expandafter\def\csname cf@../Scidonia/FracPerm/StateVar/Model.lean@369-370\endcsname{4ea30c66f461}
\expandafter\def\csname cf@../Scidonia/FracPerm/StateVar/Model.lean@423-424\endcsname{583a4eebbeaa}
\expandafter\def\csname cf@../Scidonia/FracPerm/StateVar/Model.lean@522-523\endcsname{bd6437b6df69}
\expandafter\def\csname cf@../Scidonia/FracPerm/StateVar/Model.lean@529-530\endcsname{a0fef9ee5433}
\expandafter\def\csname cf@../Scidonia/FracPerm/StateVar/Model.lean@535-535\endcsname{57d391554422}
\expandafter\def\csname cf@../Scidonia/FracPerm/StateVar/Model.lean@550-551\endcsname{b15f7df298db}
\expandafter\def\csname cf@../Scidonia/FracPerm/StateVar/Model.lean@597-600\endcsname{ef31681ee890}
\expandafter\def\csname cf@../Scidonia/FracPerm/StateVar/Model.lean@607-608\endcsname{e53944a7355a}
\expandafter\def\csname cf@../Scidonia/FracPerm/StateVar/Model.lean@582-587\endcsname{d4cc154e74a3}
\expandafter\def\csname cf@../Scidonia/FracPerm/StateVar/Model.lean@642-643\endcsname{14055dfba24e}
\expandafter\def\csname cf@../Scidonia/FracPerm/StateVar/Model.lean@656-657\endcsname{ba995450fed5}
\expandafter\def\csname cf@../Scidonia/FracPerm/StateVar/Model.lean@667-668\endcsname{611a3eae16f6}
\expandafter\def\csname cf@../Scidonia/FracPerm/StateVar/Heap.lean@102-103\endcsname{5b740228b3d1}
\expandafter\def\csname cf@../Scidonia/FracPerm/StateVar/Heap.lean@118-119\endcsname{9d55df76baec}
\expandafter\def\csname cf@../Scidonia/FracPerm/StateVar/Heap.lean@126-127\endcsname{548b3c12e0e5}
\expandafter\def\csname cf@../Scidonia/FracPerm/StateVar/Heap.lean@135-140\endcsname{f8af60b72496}
\expandafter\def\csname cf@../Scidonia/FracPerm/StateVar/Heap.lean@157-162\endcsname{8f1dc76ac072}
\expandafter\def\csname cf@../Scidonia/FracPerm/StateVar/Heap.lean@168-171\endcsname{ce7f92922125}
\expandafter\def\csname cf@../Scidonia/FracPerm/StateVar/Heap.lean@214-214\endcsname{c875ce15c57e}
\expandafter\def\csname cf@../Scidonia/FracPerm/StateVar/Array.lean@80-81\endcsname{5084cdad01e1}
\expandafter\def\csname cf@../Scidonia/FracPerm/StateVar/Array.lean@91-92\endcsname{486d61359463}
\expandafter\def\csname cf@../Scidonia/FracPerm/StateVar/Array.lean@113-114\endcsname{dda3851818c2}
\expandafter\def\csname cf@../Scidonia/FracPerm/StateVar/Array.lean@120-121\endcsname{8d41205432b3}
\expandafter\def\csname cf@../Scidonia/FracPerm/StateVar/Array.lean@165-166\endcsname{8041063aa6d3}
\expandafter\def\csname cf@../Scidonia/FracPerm/StateVar/Array.lean@194-195\endcsname{6a9ad9c69b16}
\expandafter\def\csname cf@../Scidonia/FracPerm/StateVar/Array.lean@216-220\endcsname{f98f33160e64}
\expandafter\def\csname cf@../Scidonia/FracPerm/StateVar/Universe.lean@86-87\endcsname{55bff5c33c9b}
\expandafter\def\csname cf@../Scidonia/FracPerm/StateVar/Universe.lean@103-103\endcsname{e23eb9c0c689}
\expandafter\def\csname cf@../Scidonia/FracPerm/StateVar/Universe.lean@153-154\endcsname{90e6c7c31d1e}
\expandafter\def\csname cf@../Scidonia/FracPerm/StateVar/Universe.lean@159-160\endcsname{6b6a3bfec516}
\expandafter\def\csname cf@../Scidonia/FracPerm/StateVar/Universe.lean@178-181\endcsname{e882d85a6188}
\expandafter\def\csname cf@../Scidonia/FracPerm/StateVar/Universe.lean@185-187\endcsname{e28b3a42614c}
\expandafter\def\csname cf@../Scidonia/FracPerm/StateVar/Universe.lean@193-199\endcsname{1fb0ee566945}
\expandafter\def\csname cf@../Scidonia/FracPerm/StateVar/Universe.lean@273-277\endcsname{960c08503b82}
\expandafter\def\csname cf@../Scidonia/FracPerm/StateVar/Slot.lean@67-67\endcsname{191b5865ef60}
\expandafter\def\csname cf@../Scidonia/FracPerm/StateVar/Slot.lean@73-73\endcsname{2e5a81e5133f}
\expandafter\def\csname cf@../Scidonia/FracPerm/StateVar/Slot.lean@82-82\endcsname{d66728bab821}
\expandafter\def\csname cf@../Scidonia/FracPerm/StateVar/Slot.lean@131-132\endcsname{777bd7ea6843}
\expandafter\def\csname cf@../Scidonia/FracPerm/StateVar/Slot.lean@152-156\endcsname{6f3e387a51d5}
\expandafter\def\csname cf@../Scidonia/FracPerm/StateVar/Slot.lean@169-174\endcsname{dca6a2e07dbf}
\expandafter\def\csname cf@../Scidonia/FracPerm/StateVar/Slot.lean@192-192\endcsname{8b3813d52365}
\expandafter\def\csname cf@../Scidonia/FracPerm/StateVar/Layout5.lean@190-191\endcsname{531d56988e55}
\expandafter\def\csname cf@../Scidonia/FracPerm/StateVar/Layout5.lean@194-195\endcsname{81d23c37433c}
\expandafter\def\csname cf@../Scidonia/FracPerm/StateVar/Layout5.lean@198-199\endcsname{3d4fd270215b}
\expandafter\def\csname cf@../Scidonia/FracPerm/StateVar/Layout5.lean@210-211\endcsname{3b81713f1639}
\expandafter\def\csname cf@../Scidonia/FracPerm/StateVar/Layout5.lean@221-222\endcsname{9381cefb1a86}
\expandafter\def\csname cf@../Scidonia/FracPerm/StateVar/Layout5.lean@270-271\endcsname{12f178d518f8}
\expandafter\def\csname cf@../Scidonia/FracPerm/StateVar/Layout5.lean@474-475\endcsname{2e37ef6f1c71}
\expandafter\def\csname cf@../Scidonia/FracPerm/StateVar/Layout5.lean@511-512\endcsname{851dd834f567}
\expandafter\def\csname cf@../Scidonia/FracPerm/StateVar/Layout5.lean@1027-1028\endcsname{0b9db9893a57}
\expandafter\def\csname cf@../Scidonia/FracPerm/StateVar/Layout5.lean@1055-1056\endcsname{79dee75ddae1}
\expandafter\def\csname cf@../Scidonia/FracPerm/StateVar/Layout5.lean@1243-1246\endcsname{c68bb8fcb258}
\expandafter\def\csname cf@../Scidonia/FracPerm/StateVar/Layout5.lean@1265-1267\endcsname{75d40ce95c38}
\expandafter\def\csname cf@../Scidonia/FracPerm/StateVar/Layout5.lean@1283-1284\endcsname{030cb2e59e4d}
\expandafter\def\csname cf@../Scidonia/FracPerm/StateVar/Layout5.lean@1361-1364\endcsname{0f90aede78e5}
\expandafter\def\csname cf@../Scidonia/FracPerm/StateVar/Layout5.lean@1379-1382\endcsname{7501cb075a09}
\expandafter\def\csname cf@../Scidonia/FracPerm/StateVar/Layout5.lean@1402-1405\endcsname{7f8964e9f0e6}
\expandafter\def\csname cf@../Scidonia/FracPerm/StateVar/Sealed.lean@125-127\endcsname{a150e108c8bc}
\expandafter\def\csname cf@../Scidonia/FracPerm/StateVar/Sealed.lean@132-135\endcsname{fa3821e2108b}
\expandafter\def\csname cf@../Scidonia/FracPerm/StateVar/Sealed.lean@152-153\endcsname{232c9c093968}
\expandafter\def\csname cf@../Scidonia/FracPerm/StateVar/Sealed.lean@157-158\endcsname{50d3a390ba9d}
\expandafter\def\csname cf@../Scidonia/FracPerm/StateVar/Sealed.lean@191-193\endcsname{865f8eda760b}
\expandafter\def\csname cf@../Scidonia/FracPerm/StateVar/Sealed.lean@199-202\endcsname{0af04bcb3c10}
\expandafter\def\csname cf@../Scidonia/FracPerm/StateVar/Sealed.lean@213-217\endcsname{4b2ed661cd73}
\expandafter\def\csname cf@../Scidonia/FracPerm/StateVar/Sealed.lean@234-237\endcsname{9b83be61cae0}
\expandafter\def\csname cf@../Scidonia/FracPerm/StateVar/Sealed.lean@251-253\endcsname{e1e96f07c3d9}
\expandafter\def\csname cf@../Scidonia/FracPerm/StateVar/RefCount.lean@740-741\endcsname{c2bd52ed1e21}
\expandafter\def\csname cf@../Scidonia/FracPerm/StateVar/RefCount.lean@765-765\endcsname{c1c6e144a166}
\expandafter\def\csname cf@../Scidonia/FracPerm/StateVar/RefCount.lean@771-771\endcsname{d9b5c6bdb52b}
\expandafter\def\csname cf@../Scidonia/FracPerm/StateVar/RefCount.lean@779-779\endcsname{fbe0489cba29}
\expandafter\def\csname cf@../Scidonia/FracPerm/StateVar/RefCount.lean@787-788\endcsname{163c10511c27}
\expandafter\def\csname cf@../Scidonia/FracPerm/StateVar/RefCount.lean@804-805\endcsname{aa9dea32f52b}
\expandafter\def\csname cf@../Scidonia/FracPerm/StateVar/RefCount.lean@819-823\endcsname{3352aac86792}
\expandafter\def\csname cf@../Scidonia/FracPerm/Structured/Model.lean@292-293\endcsname{331b6c9d7cb0}
\expandafter\def\csname cf@../Scidonia/FracPerm/Structured/Model.lean@297-298\endcsname{51d76a5567c4}
\expandafter\def\csname cf@../Scidonia/FracPerm/Structured/Model.lean@203-204\endcsname{b6db30baf018}
\expandafter\def\csname cf@../Scidonia/FracPerm/Structured/Model.lean@365-366\endcsname{8d3f6d606d28}
\expandafter\def\csname cf@../Scidonia/FracPerm/Structured/Model.lean@372-373\endcsname{0df1d9c5942f}
\expandafter\def\csname cf@../Scidonia/FracPerm/Structured/Model.lean@415-415\endcsname{4210bba08da7}
\expandafter\def\csname cf@../Scidonia/FracPerm/Structured/Model.lean@430-431\endcsname{469cb37dc3cf}
\expandafter\def\csname cf@../Scidonia/FracPerm/Structured/Model.lean@549-552\endcsname{5aa96b675f45}
\expandafter\def\csname cf@../Scidonia/FracPerm/Structured/Model.lean@559-560\endcsname{08b6b25ac8b0}
\expandafter\def\csname cf@../Scidonia/FracPerm/Structured/Model.lean@463-468\endcsname{4f499282495f}
\expandafter\def\csname cf@../Scidonia/FracPerm/Structured/Model.lean@565-567\endcsname{f03a2399da1a}
\expandafter\def\csname cf@../Scidonia/FracPerm/Structured/Model.lean@572-573\endcsname{3908df01cf88}
\expandafter\def\csname cf@../Scidonia/FracPerm/Structured/Model.lean@586-587\endcsname{9ff29af667c0}
\expandafter\def\csname cf@../ScidoniaTest/FracPerm.lean@96-96\endcsname{26c4a5c56142}
\expandafter\def\csname cf@../ScidoniaTest/FracPerm.lean@110-110\endcsname{c7072cc01773}
\expandafter\def\csname cf@../ScidoniaTest/FracPerm.lean@124-124\endcsname{667c4844b79c}
\expandafter\def\csname cf@../ScidoniaTest/FracPerm.lean@136-136\endcsname{af579ae4aafd}
\expandafter\def\csname cf@../ScidoniaTest/FracPerm.lean@149-149\endcsname{7cca48d0769a}
\expandafter\def\csname cf@../ScidoniaTest/FracPerm.lean@160-161\endcsname{58d6bb9503f6}
\expandafter\def\csname cf@../ScidoniaTest/FracPerm.lean@313-313\endcsname{84be31499770}
\expandafter\def\csname cf@../ScidoniaTest/StateVar.lean@281-281\endcsname{846203002cf7}
\expandafter\def\csname cf@../ScidoniaTest/StateVar.lean@287-287\endcsname{6247b6d47e02}
\expandafter\def\csname cf@../ScidoniaTest/StateVar.lean@293-293\endcsname{8ad7a3493477}
\expandafter\def\csname cf@../ScidoniaTest/Model5.lean@335-336\endcsname{3fb4624b4fe2}
\expandafter\def\csname cf@../ScidoniaTest/Model5.lean@353-353\endcsname{a10fbeeae676}
\expandafter\def\csname cf@../ScidoniaTest/Count.lean@44-44\endcsname{743326bf0f37}
\expandafter\def\csname cf@../ScidoniaTest/Count.lean@56-57\endcsname{ce08ec48a432}
\expandafter\def\csname cf@../ScidoniaTest/Count.lean@72-73\endcsname{6ef5df07669f}
\expandafter\def\csname cf@../ScidoniaTest/Count.lean@125-125\endcsname{a34eea8628ee}
\expandafter\def\csname cf@../ScidoniaTest/Count.lean@133-135\endcsname{018ed4c58639}
\expandafter\def\csname cf@../Scidonia/FracPerm/StateVar/Description.lean@24-46\endcsname{952632ccc0d9}
\expandafter\def\csname cf@../Scidonia/FracPerm/StateVar/Layout5.lean@91-119\endcsname{07c4092830a2}
\expandafter\def\csname cf@../Scidonia/FracPerm/StateVar/Layout5.lean@122-138\endcsname{4a24d4844560}
\expandafter\def\csname cf@../Scidonia/FracPerm/StateVar/Layout5.lean@147-166\endcsname{4f032786efd0}
\expandafter\def\csname cf@../Scidonia/FracPerm/StateVar/Layout5.lean@169-177\endcsname{0feaed59d5eb}
\expandafter\def\csname cf@../Scidonia/FracPerm/Ledger.lean@16-22\endcsname{4eb8cea25657}
\expandafter\def\csname cf@../Scidonia/FracPerm/Hoare.lean@107-132\endcsname{0303bf6cccb7}
\expandafter\def\csname cf@../Scidonia/FracPerm/Hoare.lean@145-152\endcsname{a45fb2b10d38}
\expandafter\def\csname cf@../Scidonia/FracPerm/Hoare.lean@178-184\endcsname{aef4c8295edc}
\expandafter\def\csname cf@../Scidonia/FracPerm/Hoare.lean@229-266\endcsname{8cc85285784b}
\expandafter\def\csname cf@../Scidonia/FracPerm/Counter.lean@22-49\endcsname{079f0ad3e32d}
\expandafter\def\csname cf@../Scidonia/FracPerm/StateVar/Value.lean@46-67\endcsname{bfb6d10b4c7b}
\expandafter\def\csname cf@../Scidonia/FracPerm/StateVar/Value.lean@68-82\endcsname{1f6a83f938ee}
\expandafter\def\csname cf@../Scidonia/FracPerm/StateVar/Class.lean@59-91\endcsname{3128a5f31a98}
\expandafter\def\csname cf@../Scidonia/FracPerm/StateVar/Heap.lean@44-47\endcsname{63b5d37b4f3c}
\expandafter\def\csname cf@../Scidonia/FracPerm/StateVar/Heap.lean@49-68\endcsname{0b25f74df449}
\expandafter\def\csname cf@../Scidonia/FracPerm/StateVar/Heap.lean@70-77\endcsname{6c7df423d965}
\expandafter\def\csname cf@../Scidonia/FracPerm/StateVar/Heap.lean@79-89\endcsname{1f6a8050c45e}
\expandafter\def\csname cf@../Scidonia/FracPerm/StateVar/Surface.lean@38-59\endcsname{2ff677e38d0a}
\expandafter\def\csname cf@../Scidonia/FracPerm/StateVar/Array.lean@57-76\endcsname{4d421a7cc316}
\expandafter\def\csname cf@../Scidonia/FracPerm/StateVar/Universe.lean@111-117\endcsname{55053fa59189}
\expandafter\def\csname cf@../Scidonia/FracPerm/StateVar/Universe.lean@210-240\endcsname{e1a2ef0a6af5}
\expandafter\def\csname cf@../Scidonia/FracPerm/StateVar/Slot.lean@46-63\endcsname{d3e2d380990e}
\expandafter\def\csname cf@../Scidonia/FracPerm/StateVar/Sealed.lean@52-60\endcsname{713c79572283}
\expandafter\def\csname cf@../Scidonia/FracPerm/StateVar/Sealed.lean@66-74\endcsname{9ec521f46635}
\expandafter\def\csname cf@../Scidonia/FracPerm/StateVar/Sealed.lean@83-90\endcsname{de0a942c6e6d}
\expandafter\def\csname cf@../Scidonia/FracPerm/StateVar/Layout5.lean@203-208\endcsname{85a1cf2fd656}
\expandafter\def\csname cf@../Scidonia/FracPerm/StateVar/Layout5.lean@298-308\endcsname{5af8e2fc8c99}
\expandafter\def\csname cf@../Scidonia/FracPerm/StateVar/Layout5.lean@457-471\endcsname{118d1cbd1fc1}
\expandafter\def\csname cf@../Scidonia/FracPerm/StateVar/Layout5.lean@539-555\endcsname{d010bd5a10fa}
\expandafter\def\csname cf@../Scidonia/FracPerm/StateVar/Layout5.lean@679-684\endcsname{61f9875d02fa}
\expandafter\def\csname cf@../Scidonia/FracPerm/StateVar/Layout5.lean@710-729\endcsname{95f71547697b}
\expandafter\def\csname cf@../Scidonia/FracPerm/StateVar/Layout5.lean@772-784\endcsname{20bf19049e08}
\expandafter\def\csname cf@../Scidonia/FracPerm/StateVar/RefCount.lean@713-728\endcsname{9729f794a150}
\expandafter\def\csname cf@../Scidonia/FracPerm/Structured/Surface.lean@127-149\endcsname{c9cdc256f152}
\expandafter\def\csname cf@../Scidonia/FracPerm/Structured/Surface.lean@166-179\endcsname{f8b42168ef1a}
\expandafter\def\csname cf@../Scidonia/FracPerm/Structured/Surface.lean@263-270\endcsname{dc0b62954321}

\def\codeframe[linerange=#1-#2]#3{%
  \edef\cfkey{cf@#3@#1-#2}%
  \ifcsname\cfkey\endcsname\else
    \PackageError{paper}{no generated frame for #3:#1-#2}%
      {Run python3 gen-frames.py in Paper/.}\fi
  \edef\cffile{frames/\csname\cfkey\endcsname.lean}%
  \par\smallskip\noindent
  \begin{minipage}{\linewidth}%
    \inputminted{lean}{\cffile}%
  \end{minipage}\par}
% A path is set with \path so that it may break at its slashes: a long title and
% a long module path on one caption line would otherwise overrun the margin.
\usepackage{url}\urlstyle{same}
% A body frame is labelled with the file and lines it quotes:
\newcommand{\sourceline}[3]{\par\smallskip\noindent\textbf{#1}\hfill
  {\footnotesize\itshape\path{#2}, lines #3}\par\nobreak}
\usepackage[colorlinks=true,linkcolor=blue!45!black,citecolor=blue!45!black,urlcolor=blue!55!black]{hyperref}
\usepackage{cleveref}

\theoremstyle{definition}
\newtheorem{definition}{Definition}[section]
\newtheorem{example}[definition]{Example}
\newtheorem{remark}[definition]{Remark}
\theoremstyle{plain}
\newtheorem{theorem}[definition]{Theorem}
\newtheorem{lemma}[definition]{Lemma}
\newtheorem{proposition}[definition]{Proposition}
\newtheorem{corollary}[definition]{Corollary}

\newcommand{\lrc}{\textup{$\lambda$RC}}
\newcommand{\lpure}{\textup{$\lambda$Pure}}
\newcommand{\Loc}{\mathsf{Loc}}
\newcommand{\Obj}{\mathsf{Obj}}
\newcommand{\rc}{\mathrm{rc}}
\newcommand{\held}{\mathsf{held}}
\newcommand{\Val}{\mathsf{Val}}
\newcommand{\ers}[1]{\lfloor #1 \rfloor}
\newcommand{\cm}{\mathrel{\approx}}
\newcommand{\step}{\longrightarrow}
\newcommand{\stepstar}{\twoheadrightarrow}
\newcommand{\Lean}[1]{\textsf{#1}}
\newcommand{\code}[1]{\texttt{#1}}
\newcommand{\Code}[1]{\texttt{#1}}
\newcommand{\doi}[1]{\textsc{doi}:\,#1}

\title{Unobservable effects\\[2pt]
\large Fractional permissions as a discipline for in-place update in a pure language}
\author{Gavin Mendel-Gleason\\\small Scidonia}
\date{2026}

\begin{document}
\maketitle

\begin{abstract}
A pure functional language can update in place. Lean's \code{Array.set!} denotes a
new array, and its compiled code stores into the existing buffer whenever that
buffer has exactly one reference. Both statements are true at once because the
store is \emph{unobservable}: nothing in the source language can distinguish a
value that was mutated in place from one that was copied, provided the old
value had no other reference. That condition, however, is a fact about the
compiled code rather than about the meaning of the program, and the compiler
decides it per site, from use counts, aliasing and projections. This paper
presents a discipline that makes the condition hold where a programmer wants it,
and that is erased before it runs. The discipline is a fractional permission
protocol in a ghost ledger, packaged in a Hoare monad whose programs can only be
built from its own primitives and runs only through one entry point. It is
stated once as a fact about compilation to a reference-counted intermediate
language, which is where the count is visible, and then lifted: over an
arbitrary object, over an arbitrary state type whose laws are proof fields, and,
through a description universe in the sense of levitation, over a family of
shapes for which one instance suffices. A sealed class in which the library owns
the reading, the writing and the encoding of the state closes the last route by
which a client can hold an alias. Throughout, the paper separates what is proved
in Lean, what is a property of the exported surface, what is measured by running
compiled code, and what is trusted of the compiler and the run time.
\end{abstract}

\section{Introduction}

Lean's values are immutable. A term such as \code{a.set! i v} denotes a new
array: same length, same elements except at \code{i}. No term of the language can
change a value that already exists, and no program can tell whether some other
array happens to share a buffer with this one. The source language is pure, and
the compiler is free to make it fast.

The compiler is fast in a specific way. When the array that \code{set!} is
applied to has exactly one reference, the emitted code stores into that array's
buffer and returns a pointer to the same buffer. When the array has two or more
references, the emitted code allocates a copy and stores into that. The two
implementations denote the same value, so nothing has to be decided about which
one to use for the program to be correct. The only thing at stake is allocation,
and therefore speed. This is the reference-counting discipline of Lean's
runtime, described by Ullrich and de Moura \cite{ullrich2019counting}, with the
reuse pass of Perceus \cite{reinking2021perceus}: an exact count lets an
implementation perform a destructive update precisely when the updated value is
about to die.

Both statements --- the source denotes a new array, the compiled code writes ---
are therefore true, and they are compatible because the write is
\emph{unobservable}. The observation language is the pure term language, and no
term of it can name an address. Unobservability is not a matter of intent,
though; it is exactly the condition that no other reference to the buffer is
live at the moment of the store. The compiler does not check whether a
programmer \emph{wanted} a mutation. It reads the uses of the value off the code
it is compiling, and it stores in place when its count is 1.

\subsection{The problem}

That puts the programmer in an awkward position. Whether a particular write
updates in place depends on facts about the \emph{compiled} code: how many times
a variable is used in an expression, whether a projection is held across a call,
whether an alias is still live, whether a value was produced at run time or is a
compile-time constant. None of these are visible in the meaning of the program,
and none of them can be inferred from its type. A program can be correct and
slow, and a small edit that looks semantically neutral --- returning
\code{a.size} alongside \code{a.set! i v} --- can cost a full copy of the buffer.
We measured this: on the toolchain used here, \code{(a.set! i v, a.size)} copies,
while \code{a.set! i v} alone does not.

The engineering goal is therefore not to make the compiler mutate, which is not
in our gift, but to arrange programs so that the condition under which it
mutates holds where we want the mutation, and to know that it holds.

\subsection{The idea}

We give the condition a type discipline.

A write needs exclusive access to the object. Reads may be shared. That is the
shape of a \emph{fractional permission} \cite{boyland2003checking}: the whole
permission, $1$, is required to write, while any positive share may be lent to a
reader; shares split and recombine, and once every read has been returned the
shares sum to $1$ and writing is permitted again.

Lean has no linear types. A token value can be copied like any other value, so
the permission cannot live in the token. It lives instead in a \emph{ghost
ledger} of entries $(\mathit{id}, q)$ that appears only in the types and proofs
of the program, and that the compiler erases. A token is a name for an entry; a
copied token is a second name for one permission, not a second permission. The
program type is a Hoare monad whose precondition and postcondition speak about
the ledger and about the state, and the only entry point demands a proof that
the program ends holding the whole object.

The discipline is stated once, in a form that does not depend on the object: an
object supplies its read and write operations, and the permission rules are
uniform over every object. Then it is lifted in three steps.

\begin{enumerate}[leftmargin=1.4em,itemsep=2pt]
\item \emph{Any object.} An object joins the interface by naming its read and
write operations and by supplying a check that the runtime runs before an
update. The check is the identity in the logic, so no specification or proof
mentions it, and its compiled form is the exclusivity test.
\item \emph{Any state type.} A second interface makes the logical part of the
rules into proof fields of a class. An instance must prove that its initial
states are admissible, that writes preserve admissibility, that its encoding of
read results is injective, and that a read of an admissible state names no
protected location. Every theorem of the generic model is then stated once, for
every instance.
\item \emph{Any shape.} The last step is the one that answers ``can this be done
for shapes we have not written yet?''. It can, by \emph{levitation}
\cite{chapman2010levitation}: describe the shapes in a universe of
descriptions, interpret a description as a type, and prove the laws once for the
whole universe rather than once per shape. A user's type enters by a bijection
to a described type, along which the one instance's operations are transported.
When the relation to the described type is an equality rather than a bijection,
something stronger becomes possible: the library can own the reading, the
writing and the encoding of the state, and a client can then neither introduce
an alias nor choose an encoding that hides one.
\end{enumerate}

The chain of the paper is therefore: a pure language whose implementation has
\emph{unobservable} effects; a proof that the effects are unobservable,
conditional on a reference count of one; a type discipline that establishes the
condition; generalisation of that discipline from one array to an arbitrary
object, to an arbitrary state type, and to a family of shapes by levitation; and
a sealed surface where the library, rather than the client, decides what a read
means.

\subsection{Contributions}

\begin{enumerate}[leftmargin=1.4em,itemsep=2pt]
\item A rendering of the reference-counted write as a rule of a small
intermediate language \lrc, and a proof that the rule is equivalent to the pure
copy exactly when the count is one (\cref{sec:lambda}). Unobservability is thus
made into a statement with a premise, and the premise is the thing a discipline
has to deliver.
\item A permission discipline with no linear types in the host language
(\cref{sec:perm}): fractional permissions in an erased ghost ledger, a Hoare
monad sealed by private constructors and a private capability, and an
identification of what it does and does not establish --- authority is not
reference.
\item Generalisation over objects and state types (\cref{sec:object,sec:lawful}),
with the obligations that are proof fields and the negative controls that show
they are not vacuous.
\item The levitation step (\cref{sec:levitation,sec:sealed}): one lawful instance
for every description of a small universe, transport along a bijection for a
user's type, and a sealed class that removes the author's choice of read,
write and encoding altogether.
\item A worked path calculus over layouts, with take-out writes, and two
exclusivity checks that are the difference between an update in place and a
copy (\cref{sec:sealed}).
\item An explicit statement of the correspondence chain that connects the
discipline to the compiled program, with each link labelled as proved,
structural, measured or trusted (\cref{sec:evidence}).
\end{enumerate}

All of the code discussed here is Lean 4, on a pinned toolchain, in a library
that builds and whose checks are theorems rather than test assertions. The
paper quotes the library's declarations verbatim in \cref{app:listings}: 37
panels for the types and the operations, and the statements of 110 of the
library's 184 public theorems. Every frame is a line range into the library,
re-checked by \code{check-listings.py}, so the appendix cannot drift from the
source it quotes.

\subsection{How much is claimed}

The paper is careful about the status of each statement, and says so where the
statement is made. Four labels are used throughout.

\begin{itemize}[leftmargin=1.4em,itemsep=2pt]
\item \emph{Proved}: a theorem of the Lean development. For example, that a
returned read gives back exactly its fraction, or that a read of a valid state
contains no reference.
\item \emph{Structural}: a property of the exported surface, shown by probes
from another module that the elaborator refuses. For example, that no public
route returns the state a session owns.
\item \emph{Measured}: a fact about compiled code on this toolchain, established
by running an executable that observes addresses. For example, that a write on a
uniquely owned array keeps the buffer's address.
\item \emph{Trusted}: what remains outside all of it --- the compiler's
translation, the runtime's primitives, and the arithmetic of the reference
count.
\end{itemize}

No theorem here certifies that a compiled write updates in place. What is proved
is the permission protocol and the shape of the interfaces; what is measured is
the in-place behaviour of the compiled code; what connects them is a
correspondence chain that \cref{sec:evidence} states in full, link by link, with
the links that are assumptions named as assumptions.

\section{Setting: counts, and why the condition is not a typing fact}\label{sec:setting}

\subsection{A value, an address, and a count}

A Lean value of an immutable type is a fully determined piece of data. Arrays
and structures of scalar data are represented at run time as heap objects whose
fields are values, and a value that occurs several times in the program may be
represented by one object with several references. The runtime keeps an exact
count of those references. The count is not part of the value; it is part of the
machine, and the source language has no term that can read it.

The count $n$ of an object supports exactly three questions.

\begin{itemize}[leftmargin=1.4em,itemsep=2pt]
\item $n = 1$: this reference is the only one. If the code that holds it
replaces the object by a new value, nothing can still be holding the old one, so
the old object's storage may be overwritten in place. The runtime's
\code{lean\_is\_exclusive} asks this question.
\item $n \geq 2$: something else may still read the object. A write must
allocate and copy.
\item $n = 0$: the object is a compile-time constant, held in a read-only
segment. It can never be exclusive, and a write to it always copies.
\end{itemize}

A fourth case exists in the runtime and matters for anyone writing a test from
the count: a count that has overflowed while single-threaded is parked at a
sticky value and reads as ``unshared'' (\code{lean\_is\_unshared}, which is
$\mathit{rc} > 0 \lor \mathit{rc} \leq \mathsf{LEAN\_RC\_STUCK\_ST}$) while
standing for more references than it can still track. \code{isUnshared} is
strictly weaker than \code{IsExclusive}, and a write that tested the former
would lose those references. \Cref{app:listings} quotes both predicates from the
development, which vendors the Lean repository's own machine-checked model of
\code{lean\_inc\_ref\_n} and \code{lean\_dec\_ref}
\cite{lean4rcmodel} and adds the counter-level facts the write discipline rests
on: that an extra reference costs the in-place update, that dropping it restores
it, that a thread-shared count is never exclusive, and that a never-freed count
never becomes exclusive however the adjustments run.

\subsection{What the compiler reads}

The count is maintained by insertions the compiler makes while translating to
its intermediate representation: an increment wherever a value acquires a second
reference, a drop wherever one goes out of scope. Which counts those insertions
produce is decided by the \emph{use} of values in the code being compiled.
Three signals decide every case we have measured.

\begin{itemize}[leftmargin=1.4em,itemsep=2pt]
\item An \code{inc} before a mutation: the value is used again in the same
expression, or an alias is live. \code{(a.set! i v, a.size)} emits
\code{inc x; Array.set! ◾ x …; Array.size ◾ x} --- one extra use of the same
array, and a copy.
\item \code{isShared} before a whole-value rebuild: the reset/reuse pass rewrites
the \emph{incoming} cell of a constructor when that cell is exclusive, and
allocates when it is not. The pass matches on shape, so when two same-shaped
constructors appear in one expression the reused cell can go to the wrong one.
\item Take-out against a borrowed read: an element exchanged out of an array is
owned by the caller, while an element read by a borrowed accessor is not, and
the two differ by exactly one \code{inc}. Identical intent, opposite outcome.
\end{itemize}

None of this is visible in the type of the program or in its meaning as a pure
term. It is a local, syntactic property of the code the compiler emits. A source
language cannot make the property true by typing alone; it can only arrange that
the code it accepts is the code whose uses are forgiving.

We take from this a precise formulation of the goal. Give the source language a
discipline under which a well-disciplined program compiles to code in which the
object at each mutation site has no other reference. The discipline must be
checked where the program is written, must cost nothing at run time, and must
not depend on the shape of the object.

\section{\texorpdfstring{\lpure}{Pure} and \texorpdfstring{\lrc}{RC}: the store, and when it is invisible}\label{sec:lambda}

This section fixes the language in which the claim of unobservability is
stated. \lpure\ is the pure source language, and its meaning is what a client is
entitled to observe. \lrc\ is the language the compiler targets, with addresses
and reference counts and a store instruction. Erasing the addresses of a \lrc\
configuration recovers a \lpure\ value, and a store step preserves that erasure.
The condition under which it does is precisely that the stored-into object has
one reference, which is the premise a discipline must deliver.

\subsection{The pure language}

\begin{definition}[\lpure, syntax]
Types and terms are given by
\[
\begin{array}{rcl}
\tau & ::= & \mathbf{1} \mid \mathbb{N} \mid \mathbb{B} \mid \tau \times \tau \mid \tau + \tau \mid [\tau] \\[2pt]
e & ::= & x \mid c \mid \langle\rangle \mid \langle e, e \rangle \mid \mathsf{inl}\, e \mid \mathsf{inr}\, e \mid \mathsf{case}\ e\ \mathsf{of}\ \{\dots\} \\
 & & \mid \lambda x{:}\tau.\, e \mid e\,e \mid \mathsf{arr}(e_1,\dots,e_k) \mid e[e] \mid {|e|} \mid \mathsf{let}\ x = e\ \mathsf{in}\ e
\end{array}
\]
The type $[\tau]$ is that of arrays, written \code{Array τ} in the development.
Values are the customary closed forms: constants, the unit, pairs, injections,
abstractions, and arrays, where an array value $a$ is a finite map from indices
to values.
\end{definition}

\begin{definition}[the pure write]\label{def:set}
If $a$ is an array value, $i$ an index and $v$ a value, then
$a[i \leftarrow v]$ is the array agreeing with $a$ at every index except $i$,
where it is $v$, if $i$ is within bounds; and $a$ itself otherwise. There is no
case in which $a$ changes.
\end{definition}

Evaluation is the usual call-by-value big-step relation $e \Downarrow v$ for
closed terms, with \cref{def:set} as the meaning of the store syntax. Every
closed well-typed term evaluates to a value; there is no store, no allocation
and no failure in the language. It follows that the observation a client is
entitled to make of a run is a value of the pure language, and two runs that
produce the same value are indistinguishable to every client context.

\subsection{The counted language}

\lrc\ has the same terms, and additionally addresses and two bookkeeping
instructions.

\begin{definition}[\lrc, syntax and store]\label{def:lrc}
Values are extended by addresses $\ell \in \Loc$. A store $\sigma$ is a finite
map from addresses to objects, where an object is a box $\langle v_1,\dots,v_k
\rangle$ of values together with a count $n \in \mathbb{Z}$. A configuration is
a pair $\langle \sigma, e \rangle$ with $e$ a closed term. The two additions to
the syntax are $\mathsf{inc}\ \ell$ and $\mathsf{dec}\ \ell$; $\mathsf{arr}(v_1,
\dots, v_k)$ is now an allocating form.
\end{definition}

The step relation is call-by-value with evaluation contexts $C$, and the
following rules, omitting the congruence rules.

\begin{mathpar}
\inferrule*[Right=\textup{\textsc{alloc}}]
  { \ell \notin \mathrm{dom}(\sigma) }
  { \langle \sigma, C[\mathsf{arr}(v_1,\dots,v_k)] \rangle \step
    \langle \sigma[\ell \mapsto (\langle v_1,\dots,v_k \rangle, 1)], C[\ell] \rangle }

\inferrule*[Right=\textup{\textsc{inc}}]
  { \sigma(\ell) = (b, n) }
  { \langle \sigma, C[\mathsf{inc}\ \ell] \rangle \step
    \langle \sigma[\ell \mapsto (b, n{+}1)], C[\langle\rangle] \rangle }

\inferrule*[Right=\textup{\textsc{dec}}]
  { \sigma(\ell) = (b, n) \quad n \geq 2 }
  { \langle \sigma, C[\mathsf{dec}\ \ell] \rangle \step
    \langle \sigma[\ell \mapsto (b, n{-}1)], C[\langle\rangle] \rangle }

\inferrule*[Right=\textup{\textsc{dec-free}}]
  { \sigma(\ell) = (b, 1) }
  { \langle \sigma, C[\mathsf{dec}\ \ell] \rangle \step
    \langle \sigma \setminus \ell, C[\langle\rangle] \rangle }
\end{mathpar}

and, for the array write, two rules that differ only in the count.

\begin{mathpar}
\inferrule*[Right=\textup{\textsc{store-1}}]
  { \sigma(\ell) = (\langle v_1,\dots,v_k \rangle, 1) \quad i \leq k }
  { \langle \sigma, C[\ell[i] \leftarrow v] \rangle \step
    \langle \sigma[\ell \mapsto (\langle v_1,\dots,v_{i-1},v,v_{i+1},\dots,v_k \rangle, 1)],
    C[\ell] \rangle }

\inferrule*[Right=\textup{\textsc{store-\(n\)}}]
  { \sigma(\ell) = (\langle v_1,\dots,v_k \rangle, n) \quad n \geq 2 \quad i \leq k \quad
    \ell' \notin \mathrm{dom}(\sigma) }
  { \langle \sigma, C[\ell[i] \leftarrow v] \rangle \step
    \langle \sigma[\ell \mapsto (\langle v_1,\dots,v_k \rangle, n{-}1),
      \ell' \mapsto (\langle v_1,\dots,v_{i-1},v,v_{i+1},\dots,v_k \rangle, 1)],
    C[\ell'] \rangle }
\end{mathpar}

\Cref{def:lrc} omits a read rule. The one that matters here is the access that
\emph{retains} what it reads: an array read whose result is a box must
$\mathsf{inc}$ that box, otherwise the count would not account for the second
reference. Which accesses retain is again the compiler's decision, taken from
the uses of the result, and it is the same question as the \textup{\textsc{store}} rules'
count. We therefore state reads as reading a scalar leaf without a count change,
and note where a retaining read would be an \textup{\textsc{inc}} in the same style.

\subsection{Erasure, and the theorem about the store}

\begin{definition}[erasure]\label{def:erasure}
The erasure $\ers{\sigma}$ of a store is the partial function from addresses to
pure array values that unfolds each object's contents, replacing addresses by
the arrays they denote. A configuration $\langle \sigma, e \rangle$ erases to the
\lpure\ term obtained from $e$ by replacing each address with the value it
denotes under $\ers{\sigma}$. The unfolding terminates because the language has
no way to build cyclic data.
\end{definition}

Write $\langle \sigma, e \rangle \cm \langle \sigma', e' \rangle$ when the two
configurations have the same erasure and the same free variables. Since the only
observation a client can make is the value of a closed program, and that value
is a function of the erasure, $\cm$ implies indistinguishability: no \lpure\
context separates two $\cm$-related configurations.

\begin{lemma}[bookkeeping is invisible]\label{lem:bookkeeping}
The rules \textup{\textsc{alloc}}, \textup{\textsc{inc}}, \textup{\textsc{dec}} and \textup{\textsc{dec-free}}
preserve $\cm$.
\end{lemma}

\begin{proof}[Proof sketch]
\textup{\textsc{alloc}} introduces an address that is not yet reachable from $e$, so the
erasure of the term is unchanged; when the new address is later read, it denotes
the box that was written into it, which is the box the pure term denoted.
\textup{\textsc{inc}} and \textup{\textsc{dec}} change only counts, and the erasure of
\cref{def:erasure} reads objects' contents and never their counts.
\textup{\textsc{dec-free}} removes an object whose count was $1$; the address occurred
exactly once in the term being evaluated, in the $\mathsf{dec}$, and erasure
replaces addresses by denotations, so an address occurring nowhere contributes
nothing. Each case is by induction on the evaluation context.
\end{proof}

\begin{theorem}[the store is invisible exactly at count one]\label{thm:store}
Let $\langle \sigma, C[\ell[i] \leftarrow v] \rangle$ be a configuration, with
$\sigma(\ell) = (b, n)$ and $i$ in bounds. Write $c_1$ for the result of a
\textup{\textsc{store-1}} step and $c_2$ for the result of a \textup{\textsc{store-\(n\)}} step
taken in its place. Then $c_1 \cm c_2$ if and only if $n = 1$.
\end{theorem}

\begin{proof}
The forward direction is the interesting one. Suppose $n \geq 2$. Then there is
another occurrence of $\ell$ in the term, since the count exceeds the reference
held by the store site. That occurrence is not inside $v$, whose erasure is
fixed by $v$ alone; it is elsewhere in $C$. After \textup{\textsc{store-1}}, the address
$\ell$ still denotes the updated box at that other occurrence, while after
\textup{\textsc{store-\(n\)}} it denotes the box as it was, and the updated box lives at
the fresh address $\ell'$. The two erasures differ at the other occurrence, so
they are not $\cm$-related. For $n = 1$ there is no such other occurrence, so
the erasure after \textup{\textsc{store-1}} is the erasure of the term with the updated
box at $\ell$; which is also the erasure after \textup{\textsc{store-\(n\)}}, and
\cref{lem:bookkeeping} covers the rest.
\end{proof}

\begin{corollary}[unobservability]\label{cor:unobs}
If a program's \lrc\ execution reaches every store with count one, then its
result is a value of \lpure, and the choice between storing in place and copying
at each of those sites is invisible to every pure context.
\end{corollary}

\Cref{thm:store} says what the discipline has to achieve, and says it without
mentioning types: the count at the store, and nothing else. It also explains why
a discipline is needed at all rather than an assumption. The count is
maintained by the compiler from the uses of values; the source program's meaning
does not determine it, and no amount of care in the \emph{meaning} of the
program will move it. What a discipline can do is constrain the code so that the
uses are the forgiving ones. That is what the next section builds.

\begin{remark}
\Cref{thm:store} is a theorem about \lrc\ as defined here, not about Lean's
implementation. The implementation is the subject of \cref{sec:evidence}, where
the correspondences between this calculus, the emitted code and the runtime are
either measured or named as assumptions.
\end{remark}

\section{A permission discipline, erased}\label{sec:perm}

The discipline has to establish the premise of \cref{thm:store} without
measuring anything at run time, without a linear type system in the host
language, and without restricting which objects it applies to. This section
gives the core of it: fractions in a ghost ledger, a program type that speaks
about the ledger, and a sealing discipline that keeps a client from obtaining a
program except through the primitives.

\subsection{Fractions, held fractions, and the ledger}

\begin{definition}[permission]\label{def:frac}
A \emph{fractional permission} is a rational $q$ with $0 < q \leq 1$. A write
requires the whole permission; a read requires a positive share. Shares lend and
return additively.
\end{definition}

The reading we use is Boyland's \cite{boyland2003checking}: a permission is a
share of an object, so authorisation to write is the whole share and
authorisation to read is any part of it. What the discipline has to prove is
that the shares that are out at the moment of a write amount to nothing.

\begin{definition}[ledger and held fraction]\label{def:ledger}
A \emph{ledger} is a list of pairs $(\mathit{id}, q)$, one for each read that is
currently lent out. The \emph{held} fraction of a ledger is
\[
\held(l) \;=\; 1 - \sum_{(n,q) \in l} q .
\]
\end{definition}

Two facts about \cref{def:ledger} are what the protocol rests on: lending a
fraction lowers the held fraction by exactly that fraction, and returning an
entry raises it by exactly the same fraction. The latter is
\code{Ledger.held\_erase\_of\_mem}, quoted in \cref{app:listings}: for any
$(n,q) \in l$, $\held(l.\mathit{erase}\,(n,q)) = \held(l) + q$. The ledger is
thus a complete account of how much of the object is unavailable to a writer,
and the sum returns to its starting point once every entry is erased.

\subsection{Why the permission cannot live in the token}

The obvious design puts the permission in the value that represents it: a token
of some linear type, consumed by a release. That design does not survive contact
with Lean, and the reason shapes everything that follows.

Lean has no linear, affine or uniqueness typing. A value can always be copied.
If the permission were a token carrying the fraction, then \code{let t' := t}
would produce a second permission out of one, and nothing in the type system
would notice. This is not a limitation of our particular interface; it is a
property of the host language, and it rules out the direct translation of any
ownership discipline that assumes linear values \cite{barensen1993uniqueness,
bernardy2018linear}.

The permission therefore lives outside the values, in a place the language can
reason about but the program cannot inspect: the type of the program, and the
proofs attached to it. The token becomes a \emph{name} for a ledger entry. A
copied token is a second name for the same entry; it is usable only while the
entry is present, and a release removes the entry, so a second release of the
same token has a false precondition. There is no second permission, only a
second handle.

The same argument applies to the ghost identifier. The fraction alone would not
say which outstanding read a token names, since two concurrent reads may hold
equal fractions, so each entry carries an identifier chosen by the caller. There
is no runtime counter for it: the ledger does not exist at run time, so the
identifiers must be chosen in the proof rather than drawn from the machine.

\subsection{The program type}

\begin{definition}[pre- and postcondition]\label{def:prepost}
$\mathsf{Pre}\ \sigma = \mathsf{Ledger} \to \sigma \to \mathsf{Prop}$ and
$\mathsf{Post}\ \sigma\ \beta = \mathsf{Ledger} \to \sigma \to \beta \to
\mathsf{Ledger} \to \sigma \to \mathsf{Prop}$. Both are predicates over the
ghost ledger and the state, and both live in $\mathsf{Prop}$.
\end{definition}

\begin{definition}[program]\label{def:hst}
$\mathsf{HSt}\ \sigma\ P\ \beta\ Q$ is the type of a program over the state
$\sigma$ whose precondition is $P$ and postcondition $Q$. Its constructor is
private to the module that defines it, so a client cannot build a program out of
nothing, and its single accessor is private and demands a \emph{capability}
$\mathsf{Cap}$, a private proposition that only the module can inhabit and that
the compiler erases. The only entry point from outside is
\[
\mathsf{exec} : \mathsf{HSt}\ \sigma\ P\ \beta\ Q \to (s : \sigma) \to P\,[\,]\,s
  \to (\forall b\ l'\ s',\ Q\,[\,]\,s\ b\ l'\ s' \to \held(l') = 1) \to \beta \times \sigma .
\]
\end{definition}

Two things about \cref{def:hst} are load-bearing, and both are refinements that
were forced by a failure rather than chosen in advance.

The private constructor stops a client from \emph{building} a program. It does
not stop a client from \emph{running} one: Lean's generated eliminators for a
structure expose its fields, so a private field alone still leaves the accessor
callable. Hence the capability. \code{Cap} is an uninhabited private proposition
whose only constructor is in the defining module, so an accessor that demands
one cannot be called from outside, and the eliminators cannot supply one either.
Both refusals are pinned in the source by checks that the elaborator must reject,
so a future change that reopens one fails the build.

The precondition of \code{exec} is the empty ledger, and its closing obligation
is that the program ends holding the whole object. That is the whole of the
linearity this design has: not linearity of tokens, which Lean cannot give, but
exactness of the ledger, which it can. A program that never returns a read fails
the closing obligation; a program that returns one twice fails the precondition
of the second release; a program that reads with a token whose entry is gone
fails the precondition of the read.

\subsection{The primitives}

Every primitive carries its specification in its type. The four that matter are
the following, quoted verbatim with their bodies in \cref{app:listings}.

\begin{definition}[primitives]\label{def:prims}
With $\sigma$ an object whose instance is in scope, and $\mathit{held}$ as in
\cref{def:ledger},
{\small\[
\begin{array}{rcll}
\code{acquire}\ n\ q & : & \mathsf{HSt}\ \sigma\ (\lambda l\_.\ n \notin \mathrm{ids}(l) \land
  0 < q \land q \leq \held(l))\ (\mathsf{ReadTok}\ n\ q)\ (\lambda l\ s\_.\ l' = (n,q)::l)
\\[2pt]
\code{read}\ t\ o & : & \mathsf{HSt}\ \sigma\ (\lambda l\_.\ (n,q) \in l)\ (\mathsf{Result}\ o)\
  (\lambda l\ s\ v\ l'\ s'.\ l' = l \land s' = s \land v = \code{observe}\ o\ s)
\\[2pt]
\code{release}\ t & : & \mathsf{HSt}\ \sigma\ (\lambda l\_.\ (n,q) \in l)\ \mathbf{1}\
  (\lambda l\ s\_.\ l' = l.\mathit{erase}\,(n,q) \land s' = s)
\\[2pt]
\code{write}\ o & : & \mathsf{HSt}\ \sigma\ (\lambda l\_.\ \held(l) = 1)\ \mathbf{1}\
  (\lambda l\ s\_.\ l' = l \land s' = \code{update}\ o\ s)
\end{array}
\]}
\end{definition}

Sequence is \code{bind}, whose precondition requires the first program's
postcondition to establish the second's precondition, and whose postcondition
composes the two existentially. Weakening is \code{conseq}: a stronger
precondition and a weaker postcondition are both sound. \code{withRead} is
\code{acquire}, the body, \code{release}, and return the body's result: it is
derived from the primitives and adds no trust, and it has the property that the
fraction comes back however the body ends, so a client cannot leak it and has no
release of its own to forget.

\begin{example}[two readers and a write]\label{ex:two}
A program that lends a third and two thirds to two readers, reads with each,
returns both, and then writes:

\par\medskip\noindent
\begin{minipage}{\linewidth}
\begin{minted}[frame=none,framesep=0pt]{lean}
acquire 0 (1/3)  >>>= fun t1 =>
acquire 1 (2/3)  >>>= fun t2 =>
read t1 0        >>>= fun v  =>
release t1       >>>= fun _  =>
release t2       >>>= fun _  =>
write (0, 9)     >>>= fun _  => ret v
\end{minted}
\end{minipage}
\par\medskip
Each \code{acquire} is justified by the held fraction: of the whole, a third is
lent and then two thirds, leaving nothing, which is exactly what the write
requires. The program is accepted; a second \code{release t1} would not be.
\end{example}

\subsection{What is refused, and how the refusals are stated}

Each way of misusing the discipline is refused, and the refusal is not merely a
tactic that gave up: it is a theorem that the obligation of the offending
program is \emph{false}. Writing while a read is outstanding, returning a token
twice, reading after returning, taking out more than is held, reusing an
identifier, and never returning a read are each a proposition about
\code{HSt.pre} or about the closing condition of \code{exec} that the
development proves false.

Stating refusals this way matters for two reasons. A check whose reason could be
a failure of the elaborator says nothing when the elaborator succeeds later for
an unrelated reason; a theorem of falsity says that no proof exists, so the
program can never be accepted, on any toolchain. And the theorem names the
obligation, so a reader can see which rule of \cref{def:prims} is doing the work.

\subsection{Erasure: what the compiler drops}

The discipline is specification. Every part of it that is a type former or a
proposition is erased by the compiler, and the erasure is observable in the
intermediate representation.

\begin{itemize}[leftmargin=1.4em,itemsep=2pt]
\item The ledger, every precondition, every postcondition, every proof and the
capability \code{Cap} are propositions or type formers; they appear as $\blacksquare$
in the emitted \code{Lean.IR}.
\item A read token is a \emph{type}, so it is not erased, but it has no fields:
it compiles to a constant scalar, the same representation as \code{Unit}, and
allocates nothing.
\item The ghost identifier and the fraction are ordinary \code{Nat} and
\code{Rat} arguments. Lean erases arguments by their \emph{type} --- proofs,
types, type formers --- and not by whether the body uses them, so these are not
erased as such. Written literally they are dropped by the compiler's arity
reduction; written symbolically, as \code{r/2}, they are not evaluated either,
because a call to a reduced entry removes the argument code as dead. They survive
as unused borrowed parameters only in a definition that is generic in them, and
they are computed and boxed only when a call cannot be brought to a reduced
entry at all, as through a procedure value.
\end{itemize}

The result is that a compiled client program threads one reference-counted
object and nothing else. The permission accounting costs no allocation on the
run-time path, which is the property that makes the design worth the trouble: a
discipline that allocated to prove that no allocation happens would be
self-defeating.

\subsection{Authority is not reference}\label{sec:gap}

The discipline establishes an \emph{authority} invariant and not a
\emph{reference} invariant, and the gap between them is where the rest of the
paper lives.

$\held(l) = 1$ says that no reader holds a share. It is a statement about the
ledger. Whether the object is exclusively owned by the machine is a statement
about the object graph, and the ledger is ghost state that the object graph
knows nothing about. Three routes defeat the inference from authority to
reference.

\begin{enumerate}[leftmargin=1.4em,itemsep=2pt]
\item The caller's binding. \code{exec} takes the state from the caller. A
caller that still holds its own reference to that state has two references, and
the write copies. Measured: the shared control case in the harness.
\item A read result that shares memory. A consumer of a read may return its
fraction honestly and keep the value it received. If that value shares mutable
structure with the state, the next write copies. The array interface's read
returns an \code{Option α}, which for an array element is a scalar, so the leak
does not arise there; it arises as soon as an object's read can return a part of
the state.
\item The library's own temporaries. A take-out write returns a pair, and while
that pair is live the extracted element has two references. The check must
therefore run after the pair has been destructured, which is a fact about the
emitted code and was found by measurement, not by reading the source.
\end{enumerate}

So the permission discipline closes one half of the problem: it makes a write's
\emph{authority} condition a proof obligation rather than a hope. The other
half --- that no other reference to the object exists --- is a property of the
interface and of the code the interface is compiled to, and \cref{sec:object}
begins the work of making it structural.

\section{Any object, not only arrays}\label{sec:object}

Nothing in \cref{sec:perm} mentions arrays. This section makes that explicit by
letting an object join the interface, and by moving the in-place check into the
object, where it can know what a write touches.

\subsection{The object interface}

\begin{definition}[object]\label{def:fracobject}
An object $\sigma$ is an instance of a class with the following fields: a type
family \code{Result} of the values its read operations yield, a reading function
\code{observe}, a writing function \code{update}, a check
\code{check : String → WriteOp → σ → σ}, and a law
$\code{check}\ \mathit{site}\ o\ s = s$. The two operation types are class
\emph{parameters} rather than fields, because an argument whose expected type is
a projection of the class is not reducible until that projection's instance is
known, and instance search does not resolve the nested goal; with the types as
parameters, a literal such as the array's index, or the pair $(0,7)$, can be
given as an operation.
\end{definition}

The check is the identity in the logic. That is what the law says, and it is why
no specification or proof mentions it: a specification states what \code{update}
does, and \code{update} is what runs after the check returns. Its compiled form
is something else entirely. The default check is
\code{assertExclusive}, whose logical content is the identity and whose compiled
content, supplied by an \code{implemented\_by} attribute, tests the object's
exclusivity and, when the process environment asks for it, reports the site and
exits. A write therefore pays one out-of-line call, one load and one branch, and
nothing on the write path reads the environment: the switch is read once, at
initialization.

\begin{definition}[the array instance]\label{def:arrayinst}
An array joins the interface as follows: a read operation is an index and yields
\code{a[i]?}; a write operation is an index and a value and sets that element;
and the check is the array's exclusivity check, which names the array's size in
its report.
\end{definition}

\begin{example}[a second object]\label{ex:counter}
A structure of a count and a log joins the interface with a check of its own:
exclusivity of the structure, and, for the operation that appends to the log,
exclusivity of the log. A projection of a constructor takes no reference of its
own, so checking the field does not itself share it. This is what the check hook
is for: the operation is a parameter, so a check can leave alone what that
operation does not write. An abort in checked mode then means that \emph{this}
write is about to copy, not merely that some part of the object is shared.
\end{example}

\subsection{The two rules that no specification can state}

Two obligations attach to any instance, and neither is enforceable from the
specifications, because both are about how Lean's values are represented rather
than about their mathematics. The interface states them in prose, in the
documentation of the class.

\begin{enumerate}[leftmargin=1.4em,itemsep=2pt]
\item \code{observe} must not return a reference to a part of the state that some
write updates in place. A read that hands back the array field of a structure
leaves an alias behind, and the next write copies that field. Sharing values
that no write mutates in place --- a \code{Nat}, or an element's own value --- is
harmless and intended; it is the mutable parts that a read result must not name.
\item \code{update} must not keep an old part of the state alive. Writing one
field with functional update syntax rebuilds the state around the other fields;
a rebuild that copies them makes the next write copy them.
\end{enumerate}

These obligations are the subject of \cref{sec:lawful} and \cref{sec:sealed},
where they become proof fields of a class and then disappear into the library.
Before that, the discipline already has a property worth stating: leaking an
alias does not make a program wrong. Lean is pure, so a missed in-place update
costs an allocation and a copy, and never a wrong result or a race. The
discipline is a performance discipline, and its failure mode is time.

\section{A notional interface: fractions internal, one runner}\label{sec:notional}

The next step removes the fraction from the client's view without removing it
from the theory. The client writes effectful-looking code over an object it
cannot name; the fractions are what justify the interleavings the semantics
admits.

\subsection{Readers and programs}

The surface has two computation types. A \code{Reader} may read: \code{get} at
an index, \code{size}, and \code{parallel}, which pairs two readers. A
\code{Program} may read and write: \code{readPhase} runs a reader, and \code{set}
writes an element and reports whether the write was in bounds. A program is run
by \code{runFresh}, which is given a size and a fill value, creates the array
itself, runs the program, and returns the result; \code{runFresh} is an
\code{IO} action precisely because the array must be created at run time, and a
compile-time constant can never be exclusive.

Both computation types carry their command tree behind a private constructor and
a private field, and the interpreter that runs a tree against an array is
private to its module. \code{runFresh} is therefore the only runner, and no
public function of the executor accepts an array. What the types do expose ---
through generated eliminators and through \code{Reader.treeOf} and
\code{Program.treeOf} --- is a tree of reads, writes, forks and pure
continuations over observations: a recipe, with no state in it. The model's own
evaluators are public and pure; they take the array from the caller, they are
what the theorems are stated about, and they are not a route to the array a
session creates.

\subsection{The semantics, and why the fractions are still there}

The model is a small-step relation over threads, where a thread is a running
read phase. It has a step for every interleaving of a fork's branches, and that
is what makes \code{parallel} notional: the shipped executor runs the left
branch and then the right, while the model admits every order, and the theorems
say the two agree.

\begin{itemize}[leftmargin=1.4em,itemsep=2pt]
\item Shares are conserved at a fork and restored at a join, and a fork halves
its parent's share. A well-formed leaf's shares sum to the thread's share and
stay positive.
\item No write step has a source while a read phase runs, and a well-formed
reading configuration leaves the writer no authority. This is the semantic
counterpart of the ledger's $\held = 1$.
\item A join becomes a run only when both children have finished.
\item A read's result is a function of its observations alone; a write preserves
the array's size and yields \code{true} exactly when it is in bounds.
\item Every interleaving computes what the sequential denotation computes. The
harness cases that arm the in-place check for this executor exit zero, which is
the run-time counterpart of the same claim.
\end{itemize}

The fractions are specification, and more so here than in \cref{sec:perm}: on
the run-time path there is neither a token nor a ledger, and the compiled
executor threads one array and the command trees it is walking. A read phase
therefore costs no allocation for permissions; only the continuations the trees
already carry.

\subsection{What has changed about the discipline}

Sealing the state changes the role of the fractions, and it is worth being
explicit because it changes the pitch.

An object's read and write operations are only reachable through \code{Reader}
and \code{Program}, and the state those operate on is created by \code{runFresh}
and never returned. A client cannot hold a reference to the session's array,
because it never receives one, and no operation returns one. Under that
restriction, the authority that the fractions account for and the exclusivity
the runtime needs no longer diverge through the client: the \emph{types} give
the write its exclusivity --- a reader has no write, and a write exists only
outside a read phase --- and the fractions account for something else, which is
\emph{sharing}: how a fork splits authority, why a join waits for both branches,
and why every interleaving agrees with the sequential denotation.

That is a cleaner story than the one \cref{sec:perm} tells, and it also makes the
remaining gap precise. Sealing removes the client as a source of aliases. It does
not remove the library's own temporaries, it does not touch the compiler, and it
does not touch the runtime. Those are the subjects of \cref{sec:sealed} and
\cref{sec:evidence}.

\begin{remark}[why \code{parallel} is notional]\label{rem:parallel}
A real concurrent backend would need more than permissions. On this toolchain,
an array captured by a \code{Task.spawn} closure answers \code{false} to the
exclusivity test for good, even after the task has been joined: it is marked as
shared between threads, and the next write returns a copy. Getting the array
back to a single owner is work the permission discipline does not do. The
structured executor is therefore sequential, and the fork/join is a semantics
statement rather than a claim about tasks.
\end{remark}

\section{Any state type: laws as proof fields}\label{sec:lawful}

\Cref{sec:object} leaves two obligations in prose: a read must not name the
mutable state, and a write must not keep an old part of it alive. This section
turns the logical content of the first into proof fields, so that the theorems
of the model hold of any instance that proves them, and states exactly how much
of the second remains.

\subsection{A language in which a reference can be written down}

The first obstacle is that the obligation is inexpressible: ``a read must not
return a reference to the state'' is a sentence about values, but the values a
read returns are of arbitrary type, and the type does not say whether a
reference is inside.

The remedy is to fix a language for read results, small enough that a reference
is a syntactic object in it.

\begin{definition}[values and detachment]\label{def:value}
The \code{Value} universe has a constructor for each shape a read may return ---
a unit, a natural, an unsigned word, a boolean, an optional, a sum, a pair ---
and one more, $\code{.loc}\ \ell$, which names a location. \code{NoRef ℓ v}
holds of a value $v$ that contains no $\code{.loc}\ \ell$ anywhere inside it,
however deeply nested: it descends through \code{some}, \code{inl}, \code{inr}
and \code{pair}, so a location hidden beneath a compound observation is caught.
\end{definition}

\subsection{The class}

\begin{definition}[lawful state variable]\label{def:lawful}
A state variable is an instance of a class with a state type $\sigma$, an
initializer type, read and write operation types, and: the type of value a read
yields; \code{initial}; a predicate \code{Valid} naming the states the variable
may be in; \code{observe}, a pure function of an operation and a state;
\code{replace}, which performs an exclusive logical replacement and reports
whether it took effect; and \code{encodeRead}, which translates a read result
into a \code{Value}. Membership is by proof of four laws.
\[
\begin{array}{rll}
\text{(L1)} & \code{initial\_valid} & \forall i,\ \code{Valid}\ (\code{initial}\ i) \\
\text{(L2)} & \code{replace\_valid} & \forall w\ s,\ \code{Valid}\ s \to \code{Valid}\ (\code{replace}\ w\ s).1 \\
\text{(L3)} & \code{encodeRead\_injective} & \forall o\ x\ y,\ \code{encodeRead}\ o\ x = \code{encodeRead}\ o\ y \to x = y \\
\text{(L4)} & \code{observe\_detached} & \forall \ell\ o\ s,\ \code{Valid}\ s \to \code{NoRef}\ \ell\ (\code{encodeRead}\ o\ (\code{observe}\ o\ s))
\end{array}
\]
\end{definition}

Each law closes a specific hole in the generic theory.

\begin{itemize}[leftmargin=1.4em,itemsep=2pt]
\item \emph{L1} is what makes a fresh session admissible. The generic theorems
assume a well-formed configuration, and a well-formed configuration includes
\code{Valid s}. For a fixed-capacity buffer it fails for an initializer of the
wrong capacity; the development carries that counterexample.
\item \emph{L2} makes reachability an induction: with L1, every state reachable
from a fresh session is admissible, and so is every heap cell after a write. It
fails for a write that grows a fixed-capacity buffer.
\item \emph{L3} is there to keep L4 from being vacuous. Without it, encoding
every result as the unit would satisfy L4 trivially and say nothing. No generic
theorem uses L3 directly.
\item \emph{L4} is the detachment the write discipline needs: a read of an
admissible state names no protected location.
\end{itemize}

Being lawful is not inherited. An existing object instance is not automatically
a lawful state variable, and none of the historical instances is claimed as one.
Membership is a proof, and the proofs are in the check library.

\subsection{The model, and what it buys}

A second model accompanies the class: command trees for readers and programs
with their binds; a small-step relation over configurations; a denotation; a
replay function, which is the one noncomputable definition in the theory, since
comparing two operations is a logical side condition; a well-formedness
predicate for reading threads with a share for each leaf; and a safety predicate
for command trees.

Stated once for every lawful instance, this model gives the same properties the
structured interface gave for one array: share conservation, exclusion, join
discipline, write preservation, trace replay and schedule independence. On top
of it sits a single-root heap, in which \code{At} relates a location to the
state it holds, \code{ReadAt} and \code{WriteAt} describe a read and a write step
against it, reads leave it unchanged, a write changes only its own cell, and a
read--write--read sequence is a snapshot: the first read keeps the value it took
from the \emph{old} state, and the second is the written one.

The array is an instance, and two bridge theorems state that the instance's
denotation is the structured executor's: the lifted trees are the trees the
interface already exposes, and the runner's equation composes the model's own
equation with the program bridge. Thus the general theory does not replace the
structured interface; it explains why that interface's theorems hold.

\begin{example}[the laws are not vacuous]\label{ex:nonvacuous}
Each law has a countermodel, and each is a theorem that a specific instance's
obligation is false, paired with the affirmative control of an instance that
satisfies it. A read that returns $\code{.loc}\ \ell$ fails L4; a write that
leaves the admissible states fails L2; an encoding that forgets what it was
given fails L3. The development proves each.
\end{example}

\subsection{What the laws do not say}

The laws constrain the \emph{encoding} and not the object graph, and the
development carries a witness to that gap.

\begin{example}[a lawful instance whose read is the state]\label{ex:leak}
Take an array of machine words. Encode it as a nested pair of words, and the
encoding is injective and contains no $\code{.loc}$ anywhere, so L3 and L4 hold;
take \code{Valid} to be trivial and \code{replace} to be the array write, and L1
and L2 hold. All four laws hold. Yet take the read to be the identity, and the
read of the state returns the protected array itself.

The instance is legal, and it is exactly what L4's proof obligation cannot catch:
Lean's logic sees values, and an alias and a copy of an array are equal values.
\Code{NoRef} is a statement about the encoding, so an author who chooses both the
encoding and the read can choose them consistently and still hand back the state.
\end{example}

\Cref{ex:leak} is the reason the second half of this paper exists. A discipline
that lets the instance author decide what a read means has an escape hatch, and
the only way to close it is to take the decision away. There are two ways to take
it away, and they differ in how the author's type is related to a type the
library knows: by a bijection, in \cref{sec:levitation}, or by an equality, in
\cref{sec:sealed}.

\section{Any shape: a description universe}\label{sec:levitation}

The lawful-state-variable class is general over types but not over shapes: each
instance proves its own laws, and an author who wants a new state type must
discharge them all. This section removes that per-shape work with the pattern of
levitation \cite{chapman2010levitation}: describe the shapes in a universe,
interpret a description as a type, and prove the laws once for the whole
universe.

\subsection{The universe and its interpretation}

\begin{definition}[descriptions]\label{def:desc}
A description is a unit, a natural, or a sum of two descriptions, and a
description denotes a type:
\[
\mathsf{Desc} ::= \mathtt{.unit} \;\mid\; \mathtt{.nat} \;\mid\;
  \mathtt{.sum}\ (\mathit{left}\ \mathit{right} : \mathsf{Desc}),
\qquad
\begin{array}{rcl}
\mathsf{El}\ \mathtt{.unit} &=& \mathbf{1} \\
\mathsf{El}\ \mathtt{.nat} &=& \mathbb{N} \\
\mathsf{El}\ (\mathtt{.sum}\ d_1\ d_2) &=& \mathsf{El}\ d_1 + \mathsf{El}\ d_2
\end{array}
\]
Three declarations, in full: the description type, its interpretation, and the
description of a slot. \code{El} is an abbreviation rather than a definition, so
it is reducible for a client: \code{slotDesc} denotes $\mathbf{1} +
\mathbb{N}$ definitionally, and a literal injection can be given at such a type
without an ascription.
\end{definition}

%!check: inductive Desc where
\sourceline{The description universe}{Scidonia/FracPerm/StateVar/Description.lean}{24--46}
\codeframe[linerange=24-46]{../Scidonia/FracPerm/StateVar/Description.lean}

The universe is deliberately non-recursive. Its types are built from the unit,
the natural and the sum alone, so a proof over it is a finite case analysis
rather than a fixpoint. Recursion is the obvious next extension and is not
attempted here.

\begin{definition}[one instance]\label{def:universeinstance}
\code{encodeEl} encodes an element of \emph{any} description as a \code{Value},
and the facts the class needs of it hold of every description:
\code{encodeEl} is injective, and it names no location. Both are proved once, by
a structural induction on the description with one case per constructor, rather
than once per description.

Because two descriptions can denote definitionally equal types, the description
indexes a wrapper: \code{Described d} is a type whose instance argument is the
description. Then
\[
\code{describedStateVar}\ d : \mathsf{LawfulStateVar}\ (\code{Described}\ d)\ (\mathsf{El}\ d)\ \mathbf{1}\ (\mathsf{El}\ d)
\]
is a single instance serving every description, and its four laws are proved
once, at the description rather than at the type.
\end{definition}

\subsection{A user's type, by transport}

One instance for the universe is not yet an instance for a type a user wrote.
The second half of the pattern carries the instance across a bijection.

\begin{definition}[transport]\label{def:transport}
Let $\tau$ be any type with functions
$\mathit{toEl} : \tau \to \mathsf{El}\ d$ and
$\mathit{ofEl} : \mathsf{El}\ d \to \tau$ that are inverse in both directions.
Then \code{transport} yields a lawful instance at $\tau$ whose operations are the
universe instance's conjugated by the isomorphism, and whose every law field is
the universe's law after an inverse law cancels the conjugation.
\end{definition}

\begin{example}[the slot, by transport]\label{ex:slot}
A slot is a genuine inductive type --- empty, or holding a natural payload ---
described by a unit summed with a natural. Its instance is one line:

\par\medskip\noindent
\begin{minipage}{\linewidth}
\begin{minted}[frame=none,framesep=0pt]{lean}
transport slotDesc Slot encode decode decode_encode encode_decode
\end{minted}
\end{minipage}
\par\medskip
Every operation is the universe instance's conjugated by \code{encode} and
\code{decode}, and every law is the universe's, so nothing specific to the slot's
description is proved. The theorems of the generic model then apply at it: a read
after a write returns the written value; a snapshot survives, with the old value
still held by the pre-write heap and the cell holding the new one; and a single
root cell changes shape, $\mathit{empty} \to \mathit{full}\ n \to
\mathit{empty}$.

A second datatype costs only its isomorphism. The development checks a transport
for the booleans at the description $\mathtt{.sum}\ \mathtt{.unit}\
\mathtt{.unit}$, which the slot does not use, so the transport is not exercised
only by its one client.
\end{example}

This is the pattern in miniature, and it is worth naming the three roles, because
the remaining sections scale it. There is a \emph{model} --- the description
universe and its interpretation. There is a \emph{bijection} between the model
and the user's type. And there is a \emph{transfer}: observations and updates are
stated in the model and proved once there, then carried to the user's type along
the bijection, so that the operation the user calls runs on the user's own type
while the reason it is correct lives in the model.

\subsection{Where the bijection is not enough}

A transport makes a user's type lawful, but it leaves the shape of the read and
the write in the user's hands: the author supplies \code{encode}, \code{decode}
and the two inverse laws, and can still choose a read that names the state, as in
\cref{ex:leak}. What transport adds over a hand-written instance is that the laws
are not re-proved; it does not add a guarantee about what a read returns.

To get that, the relation between the author's type and the library's type has to
be tighter than a bijection. A bijection gives operations that are the model's
\emph{up to} conversion, and the conversion is the author's; an equality gives
operations that \emph{are} the model's, with the conversions erasing. Which is
what the next section builds, and the price is that the library must decide what
a read and a write mean, and therefore must have a universe in which they can be
defined for every shape at once.

\section{Owning the sense of the state}\label{sec:sealed}

This section gives the library the entire meaning of the state. It has three
parts: two universes, one for what a read may return and one for what a state
may be; a calculus of paths over the second, with reads that end in the first and
writes that take an element out before updating it; and a class whose instances
supply nothing but the layout and an equality.

\subsection{Two universes, and the asymmetry between them}

\begin{definition}[observations and layouts]\label{def:layouts}
The observation universe is closed, with one constructor per shape a read may
return and an interpretation that is total on it:
\[
\mathsf{ObsDesc} ::= \mathtt{.unit} \mid \mathtt{.nat} \mid \mathtt{.uint}
  \mid \mathtt{.bool} \mid \mathtt{.sum}\ o\ o \mid \mathtt{.option}\ o
  \mid \mathtt{.pair}\ o\ o
\]
\[
\begin{array}{rcl@{\qquad}rcl}
\mathsf{ObsEl}\ \mathtt{.unit} &=& \mathbf{1} &
\mathsf{ObsEl}\ \mathtt{.sum}\ a\ b &=& \mathsf{ObsEl}\ a + \mathsf{ObsEl}\ b \\
\mathsf{ObsEl}\ \mathtt{.nat} &=& \mathbb{N} &
\mathsf{ObsEl}\ \mathtt{.option}\ a &=& \mathsf{Option}\ (\mathsf{ObsEl}\ a) \\
\mathsf{ObsEl}\ \mathtt{.uint} &=& \mathsf{UInt64} &
\mathsf{ObsEl}\ \mathtt{.pair}\ a\ b &=& \mathsf{ObsEl}\ a \times \mathsf{ObsEl}\ b \\
\mathsf{ObsEl}\ \mathtt{.bool} &=& \mathsf{Bool}
\end{array}
\]
No constructor of it denotes an array, a structure of the state or a reference.
The layout universe has three constructors, and its interpretation is the type a
program holds:
\[
\mathsf{Layout} ::= \mathtt{.leaf}\ \mathsf{ObsDesc} \;\mid\;
  \mathtt{.struct}\ \mathsf{Layout}\ \mathsf{Layout} \;\mid\;
  \mathtt{.arr}\ \mathsf{Layout}
\]
\[
\begin{array}{rcl}
\mathsf{State}\ (\mathtt{.leaf}\ o) &=& \mathsf{ObsEl}\ o \\
\mathsf{State}\ (\mathtt{.struct}\ a\ b) &=& \mathsf{State}\ a \times \mathsf{State}\ b \\
\mathsf{State}\ (\mathtt{.arr}\ a) &=& \mathsf{Array}\ (\mathsf{State}\ a)
\end{array}
\]
A layout may contain an array; an observation may not. That asymmetry is the
whole of the next section's guarantee: an array of structures is
\code{Array (UInt64 × UInt64)} as a state and an atom as a read, so no read can
return a piece of the state.
\end{definition}

%!check: inductive ObsDesc where
\sourceline{The observation universe}{Scidonia/FracPerm/StateVar/Layout5.lean}{91--119}
\codeframe[linerange=91-119]{../Scidonia/FracPerm/StateVar/Layout5.lean}

%!check: abbrev State : Layout → Type
\sourceline{The layout universe and its interpretation}{Scidonia/FracPerm/StateVar/Layout5.lean}{122--138}
\codeframe[linerange=122-138]{../Scidonia/FracPerm/StateVar/Layout5.lean}

\subsection{Paths, and a whole-array read that cannot be written down}

\begin{definition}[path and observation path]\label{def:path}
A path into a layout is an address, indexed by the layout it walks and built by
four constructors, of which the first is the trivial path and the middle two
exist only under a product:
\[
\begin{array}{rcl}
\mathsf{Path} & : & \mathsf{Layout} \to \mathsf{Type}, \\
\mathtt{.here} & : & \mathsf{Path}\ d \\
\mathtt{.left} & : & \mathsf{Path}\ a \to \mathsf{Path}\ (\mathtt{.struct}\ a\ b) \\
\mathtt{.right} & : & \mathsf{Path}\ b \to \mathsf{Path}\ (\mathtt{.struct}\ a\ b) \\
\mathtt{.index} & : & (i : \mathbb{N}) \to \mathsf{Path}\ a \to \mathsf{Path}\ (\mathtt{.arr}\ a)
\end{array}
\]
The layout at a path's end is a function of the path alone, so it is a function
computed from the constructors rather than a search of the state:
\[
\begin{array}{rclcrcl}
\mathsf{At}\ \mathtt{.here} &=& d &\qquad&
\mathsf{At}\ (\mathtt{.left}\ p) &=& \mathsf{At}\ p \\
\mathsf{At}\ (\mathtt{.right}\ p) &=& \mathsf{At}\ p &\qquad&
\mathsf{At}\ (\mathtt{.index}\ i\ p) &=& \mathsf{At}\ p
\end{array}
\]
An \emph{observation path} is a path together with the proof that its end is a
stored observation. It is a structure of three fields, and the third is the
proof:
\[
\mathsf{ObsPath}\ d = \{\;
  \mathit{path} : \mathsf{Path}\ d,\quad
  \mathit{obs} : \mathsf{ObsDesc},\quad
  \mathit{leaf} : \mathsf{At}\ \mathit{path} = \mathtt{.leaf}\ \mathit{obs} \;\}
\]
\end{definition}

%!check: abbrev At {d : Layout} : Path d → Layout
\sourceline{Typed paths and the layout at their end}{Scidonia/FracPerm/StateVar/Layout5.lean}{147--166}
\codeframe[linerange=147-166]{../Scidonia/FracPerm/StateVar/Layout5.lean}

%!check: leaf : At path = .leaf obs
\sourceline{The observable path}{Scidonia/FracPerm/StateVar/Layout5.lean}{169--177}
\codeframe[linerange=169-177]{../Scidonia/FracPerm/StateVar/Layout5.lean}

The proof field in \cref{def:path} is what makes the interface's guarantee
structural rather than a rule. It is an equality
$\mathsf{At}\ p = \mathtt{.leaf}\ o$, and at an array or a product end
that equality has no solution, so no observation path exists there. A
whole-array read therefore cannot be \emph{written down}, let alone performed,
and the development states both sides of the same fact: a theorem that
$\mathsf{At}\ \mathtt{.here} \ne \mathtt{.leaf}\ o$ at an array root, and a
check that the elaborator rejects a use of the path that would have needed such
an equality.

The one path step whose validity depends on a particular state is the array
index, so the calculus asks separately whether a path addresses anything. That
predicate is a pure traversal: it reads the elements along the path and never
copies an array. Asking first is what lets the update itself be total, and what
lets a failed write return the state it was given.

\subsection{Reads, and the encoding the library owns}

An observation path and a state yield the library's encoding of the observation
at the path's end, or \code{none} when the path addresses nothing. The traversal
that obtains the focus is private, so the raw focus is obtained and encoded in
one step, and no function returns a state or a subtree.

The encoding has one clause per observation code, and no clause produces a
location. From the clauses alone the development proves, for every code, that an
encoded observation is injective and names no location, by an induction over the
observation universe. The consequence is the read guarantee: a successful read
names no $\code{.loc}$, so it cannot hand back a mutable part of the state. This
is the fact that \cref{ex:leak} said an author could choose away; here no author
chooses it, because the encoding is the library's.

\subsection{Writes: a specification, an implementation, and a bridge}

A write needs the state with the addressed part removed. That is the one-hole
context of the path, which for this layout universe is its derivative in
McBride's sense \cite{huet1997zipper,mcbride2001derivative}: at a leaf the hole
is the whole value; at a product one component keeps the hole and the other is
held whole; at an array the selected element is taken out and the array is held
with a placeholder, together with the index that puts it back.

\begin{definition}[the zipper as specification]\label{def:zipper}
\code{open} takes a state apart into a context and a focus, or reports the
original state for an index that addresses nothing at any depth; \code{plug}
puts a focus back. Both are private, so a client cannot obtain a raw focus.
\end{definition}

The zipper is the \emph{specification}. The \emph{implementation} is a
structurally recursive update with one clause per path constructor, and the
theorem that relates them is the bridge: on a state the private \code{open} takes
apart, the recursion is the context \code{open} reported, plugged with the focus
it reported. Everything else about writes follows from that bridge: reading back
what was written, a failed write leaving the state alone, and a write at a path
disjoint from a read path not changing what that read returns. Disjointness here
excludes the equal paths and every ancestor/descendant pair, including two paths
through the same array index, rather than merely paths whose final fields differ.

Two properties of the implementation are not about the zipper at all, and both
are decisions taken because of the compiler.

\begin{itemize}[leftmargin=1.4em,itemsep=2pt]
\item The new focus is threaded through the recursion as a \emph{value} --- the
focus of the path --- rather than as a function of the old state, so the write
path allocates no closure.
\item The recursion returns the state alone, not a state paired with the write's
report. This is what lets the compiler's reuse pass rewrite a product cell rather
than allocate a fresh one; with a pair return, the pass matched the transient
pair instead of the stored cell, and the address changed. The paper returns to
this in \cref{sec:evidence}, because it was found by measuring addresses and not
by reading the source or the emitted C.
\end{itemize}

\subsection{The array case is a take-out}

A write below an array index cannot simply read the element and write it back:
the array would still hold the element, so a write inside it would copy the
element. The element must be \emph{taken out}.

\begin{definition}[take-out]\label{def:takeout}
The array case exchanges the selected element for a placeholder, leaving a
sibling array that keeps its size and every other slot, and hands the element
back as the sole owner of it. The exchange and its inverse are private, like the
zipper: they hand back the raw element that an update below the slot works on.
\end{definition}

The array's exclusivity is not the element's. A take-out replaces one slot, and
another slot --- or a caller's own reference --- may still point at the same
element, so a write inside the element would copy it. The implementation
therefore runs \emph{two} exclusivity checks, and both are the difference between
a mutation and a copy: one on the array, before the take-out, because a shared
array is copied whole; and one on the extracted element, before the recursion
below it runs, because a shared element is what a write inside it would copy.
Each check names its site in the report.

Two remarks on the checks are worth making here, because they are the places
where the discipline meets the compiler.

First, a pair is a second reference. The exchange returns the sibling array and
the element together; while that pair is live, the element has two references,
and checking it then fires on an element that is in fact uniquely owned. The
implementation destructures the pair before anything is checked or recursed
into, so the check sees one reference. That is a fact about the emitted code, and
it was found by a harness case failing.

Second, the checks are the object interface's check from \cref{sec:object},
arrived at from the other direction: there, an instance supplied a check so that
a copy would be loud; here, the library's own recursion installs the two checks
its own writes need.

\subsection{The sealed class}

\begin{definition}[sealed state variable]\label{def:sealed}
An instance of \code{SealedStateVar} supplies two things: the layout its type is,
and a proof $\code{rep} : \sigma = \mathsf{State}\ \mathit{layout}$. It has no
field for a read, a write or an encoding, so it cannot supply one. The library
defines the operations once, from the equality: a read is the layout's read at
the transported state, and a write is the layout's write with the state cast in
and the state it returns cast out.
\end{definition}

From the equality the library also derives a lawful state variable for the sealed
type, by transporting its adapter along \code{rep}. The derived instance's read
result type and its encoding are the library's by \code{rfl}: an instance has no
field in which to put a different one, and the development proves the two
equations rather than assuming them. The read, snapshot, frame and
failure theorems are then restated for every sealed instance, each inherited by
transport rather than proved per instance.

The comparison with \cref{ex:leak} is the point of the class. The lawful
state-variable class still admits the array instance whose read is the state;
the sealed class has no field to put that read in. And the guarantee is
structural: a sealed read returns an encoding built from observations, which has
no location branch, so it cannot hand back a mutable part of the state.

\begin{example}[what the checks pin from outside]\label{ex:probes}
The property that stands in for a proof of non-escape is that the raw routes
cannot be written. From a module outside the library, the development pins, by
checks whose elaboration must fail with a pinned message: that an instance adding
an \code{observe} is rejected, that an instance adding an \code{encodeRead} is
rejected, that \code{rep} cannot name a layout the type is not, that the
\code{LawfulStateVar} fields \code{observe} and \code{encodeRead} are not names a
sealed client can even write, and that the private pieces of the path calculus
--- the context, its two directions, and the take-out and its inverse --- are not
reachable. Beside them are the affirmative controls: the class's fields are
readable, the read and the write run on both shipped instances, and the two
cast round trips hold.

Five rejections are pinned around the path calculus alone --- the context, the
open, the plug, the take-out and its inverse --- which is how the raw-focus
boundary is observed from outside the module rather than asserted.
\end{example}

The reader who wants the load-bearing detail can compare the two classes
directly. A sealed instance's \code{rep} must be an equality and not an
isomorphism, which is a real restriction and is stated as one: for a structure a
user declares, the only relation to a layout's state type is an isomorphism, and
the elaborator will not accept an equality. Such a type needs a layout that
\emph{names} it, which the current universe does not have, or a conversion on
writes whose cost has not been measured. Arrays, nested arrays, pairs and the
observation universe's leaves are usable, because the layout interpretation
denotes them definitionally; a user structure is not.

\section{Evidence, and the correspondence chain}\label{sec:evidence}

Everything so far is a design with a proof obligation attached. This section
states how the obligation is discharged in practice, where the evidence stops,
and what has to be assumed beyond it.

\subsection{What the harness measures}

The run-time claims are established by an executable that observes addresses.
For each case it records the address of the object before and after the
operation; a copy allocates a new object, so an unchanged address is evidence of
an update in place. The cases that ship are the following.

\begin{itemize}[leftmargin=1.4em,itemsep=2pt]
\item \emph{Unique writes.} The structured interface's write, and a structured
program whose read phase returns a closure over its observations, each keep their
array's address. For the layout calculus, four path writes --- a left field, a
right field, the last element and an out-of-bounds index --- keep the buffer's
address, keep the address of the product cell at the written index, and keep the
sibling struct's address; the out-of-bounds case also returns the state
unchanged.
\item \emph{Nested writes.} A write to a leaf of an inner array keeps the inner
array's own address, which is the measurement that the take-out and the
element's check are doing their work.
\item \emph{Controls that must copy.} A live second reference to the array, a
live second reference to the focused inner array, a compile-time array literal,
and a shared object in the counter example. Each asserts that the \emph{address
changes}.
\item \emph{Checked children.} Each unique case is run again in a child process
with the in-place check armed, and must exit zero having reported that the check
was enabled. The shared cases are run the same way and must \emph{terminate} the
process, naming the site. The armed control that must terminate is the
counter's.
\end{itemize}

The last group is what turns the check from an assertion into an instrument. A
unique case that exits zero with the check armed establishes that the check was
in the write path and had nothing to report; a shared case that terminates
establishes that the check fires when it should. Without the second, the first
would be consistent with a check that never runs.

Two of the harness cases exist because an earlier version passed a reading of the
source and failed at run time.

\begin{enumerate}[leftmargin=1.4em,itemsep=2pt]
\item The take-out returns a pair. While that pair is live the element has two
references, so the element's check fired on a uniquely owned element and the
whole write copied. The fix is to destructure the pair before checking or
recursing; the emitted code then releases the pair and checks the focus with one
reference.
\item With the update returning a state paired with its report, the reuse pass
rewrote the transient pair rather than the stored cell, so the stored struct was
fresh and copied --- even though the emitted C contained the exclusivity test and
the in-place field write. The fix is the state-only recursion of
\cref{sec:sealed}. The address changed while the source and the C both looked
right, which is the reason the instrument is an address assertion.
\end{enumerate}

\subsection{The shape rules}

The measurements reduce to five rules that decide whether an update takes place,
and it is worth writing them down, because four of them constrain whoever writes
the update rather than the client.

\begin{enumerate}[leftmargin=1.4em,itemsep=2pt]
\item The object was allocated at run time. A compile-time constant has count $0$
and is never exclusive, which is why the session's array is created by the
runner rather than by the caller.
\item Its count is $1$ at the mutation: no alias live across the call, no
projection of it taken and held, and no second use of the same value in the same
expression.
\item A whole-value rebuild recycles the incoming cell when that cell is
exclusive, but the pass matches on shape, so when more than one same-shaped
constructor appears in the expression the cell can go to the wrong one.
\item A write nested inside a structure needs every node on the path to be
exclusive at the moment of the write. A take-out makes the node owned; a
borrowed accessor does not.
\item The write is a fixed-size element write. A store cannot grow an array.
\end{enumerate}

Two of the five are settled by the \emph{surface} rather than by care: the write
installs an observation, so no client value --- and hence no client alias --- can
be written into the state, and there is no growth operation. That is the sense in
which the interface makes the losing shapes unwritable where types can say so.
The other three are obligations on the library's own update paths, discharged by
the shape of the recursion and checked by the harness.

It is also worth saying what a normal update looks like here, because the shape
that copies is not the shape a write takes. A write is ``check, then write'',
and the check returns a boolean and retains nothing, so the state's uses are
sequential and the last is a move: the public write compiles to a bounds check
and then the fused update, with no increment of the state at all. The copying
shape --- write and read one value at once --- cannot be expressed through the
interface, because a client has commands and not a state.

\subsection{What is diagnostic}

The in-place check is diagnostic and nothing more. It is the identity in the
logic, so no specification or proof depends on it. It reads its switch once at
process start, so a write costs a call, a load and a branch. In a run where the
switch is off it neither proves nor reports anything; and in any run its verdict
rests on the runtime's exclusivity primitive, which is part of the trusted base.
The check's purpose is to make a copy loud when someone is looking for one. The
address assertions, not the check, are what establish that an update took place.

\subsection{The chain}

The gap between the discipline and the compiled program is spanned by a chain of
statements. Each link is either proved here, a property of the exported surface,
measured on this toolchain, or assumed. Naming them individually is the point:
a single sentence such as ``the discipline gives in-place update'' would hide
four distinct assumptions.

\begin{table}[ht]
\centering
\small
\begin{tabular}{@{}p{0.05\textwidth}p{0.60\textwidth}p{0.29\textwidth}@{}}
\toprule
Link & Content & Status \\
\midrule
1 & The pure language's denotation of a program is its result & proved (\cref{sec:lambda}) \\
2 & A store with count $1$ is equivalent to the copy; with count $\ge 2$ it is not & proved (\cref{thm:store}) \\
3 & Every write site reaches its precondition, so $\held = 1$ there & proved (\cref{sec:perm}) \\
4 & $\held = 1$ at a write means no read share is outstanding & proved (the ledger and the primitives) \\
5 & No client value names the session's state or a mutable part of it & structural (\cref{sec:notional,sec:sealed}) \\
6 & The library's own temporaries add no reference at a mutation & measured, per shape (\cref{sec:evidence}) \\
7 & The compiler's counts follow the uses of the term, and its insertions implement them & assumed \\
8 & The runtime's exclusivity test is $\rc = 1$, and the store primitive stores in place exactly when it holds & assumed \\
9 & A compiled write updates in place on the shapes covered & measured (\cref{sec:evidence}) \\
\bottomrule
\end{tabular}
\end{table}

\Cref{thm:store} is what makes links 3 to 5 worth having at all. What they give
is a source program whose every write site is reached with the object
unshared by the interface --- which is link 5's content --- and with no reader
outstanding, which is links 3 and 4. Links 7 and 8 are what carry that to the
machine, and they are assumptions here rather than theorems: link 7 is the
counting discipline documented by Ullrich and de Moura
\cite{ullrich2019counting} and refined by Perceus \cite{reinking2021perceus}, and
link 8 is the runtime's contract. Link 6 is neither: it is a property of the
library's own code, and it is where the paper's measurement effort went, because
it is the only link under our control.

\begin{corollary}[discipline adequacy, conditional]\label{cor:adequacy}
Assume links 6, 7 and 8. Let $p$ be a program accepted by \code{exec} on a state
that is uniquely owned when \code{exec} is called. Then every store site in
$p$'s compiled code is reached with the object's count one, and by
\cref{cor:unobs} every effect those stores perform is unobservable to a pure
client.
\end{corollary}

\begin{remark}[what would make link 7 unconditional]\label{rem:link7}
Link 7 could be replaced by a proof against a counted calculus, with the
sealing assumptions as hypotheses: that the object never escapes, and that every
read result is built from the observation universe. That turns ``count one at a
write'' for a class of programs from measured into proved \emph{in the model},
with the correspondence to compiled executions still assumed link by link. It is
a different result from a measurement, and it should be labelled as one. The
development here does not attempt it; the vendored counter model of
\cref{app:listings} is the vocabulary such a proof would be stated in, not the
proof.
\end{remark}

\subsection{Non-claims}

A list of what is not claimed is part of the result.

\begin{itemize}[leftmargin=1.4em,itemsep=2pt]
\item No theorem here certifies that a compiled write updates in place, and none
certifies anything about Lean's object graph. The in-place evidence is a
measurement on named cases of one toolchain, and the emitted code is a
description of that toolchain rather than a specification.
\item The detachment laws are about the modelled encoding. \Code{NoRef} says that
an encoded read names no location; the sealed class removes the author's choice
of that encoding but does not prove that the library's encoding matches the
runtime's references.
\item Sealing establishes non-escape through the exported operations. It says
nothing about \code{unsafe} code, foreign functions, or values a client obtained
outside the interface.
\item The fork/join is notional. The executor is sequential, and a concurrent
backend would still face the thread-shared marking of \cref{rem:parallel}.
\item A missed in-place update costs time and never correctness. Lean is pure,
so every performance claim here is a claim about allocation.
\item The counter-level facts inherited from the vendored model inherit that
model's trust character, down to the generated certificates some of its proofs
carry; the development prints, rather than assumes, which principles each of its
own theorems rests on.
\end{itemize}

\section{Related work}\label{sec:related}

\paragraph{Fractional permissions and separation logic.}
Permissions as shares of an object are Boyland's
\cite{boyland2003checking}, in the tradition of separation logic
\cite{ohearn2001local,reynolds2002separation}. Iris \cite{jung2015iris} and its
proof-mode descendants \cite{krebbers2017essence}, VeriFast
\cite{jacobs2011verifast} and Viper \cite{muller2016viper} reason this way about
real memory. Our use of fractions is deliberately modest: no heap assertion
language, no concurrent resource algebra, and no reasoning about interference
between threads. What is borrowed is the arithmetic of shares and the
observation that a write needs the whole. What is new is where the fractions
live: in the type of a program, in a ledger the compiler erases, in a host
language with no linear types.

\paragraph{Permissions with real mutation in a proof assistant.}
F$^{\star}$'s Steel \cite{fromherz2021steel} applies fractional permissions to arrays
and extracts to C with real mutation, which makes it the closest existing system
to the first half of this paper. Verus \cite{lattuada2023verus} combines Rust's
own linear types with ghost permissions. Both are verification frameworks: the
program is proved against a specification that mentions the heap. Ours is a
discipline, not a logic: there is no specification of the object's contents, and
the only property proved is the one the writer needs --- authority, detachment,
snapshots, frames. A reader who wants to state what a program computes would
reach for a Hoare logic; the development's program type is the degenerate case
in which the specification is the permission.

\paragraph{Hoare logic for state monads.}
The program type is a Hoare state monad in the style of Swierstra
\cite{swierstra2009hoare}, who gives a Hoare logic for the state monad with
preconditions and postconditions as type indices, and of Hoare Type Theory and
Ynot \cite{nanevski2008htt,nanevski2008ynot}. Our monad is a value of that shape
with one addition: because the host language has no linear values, the
precondition is a predicate on a ghost ledger rather than a separation-logic
assertion, and the seal is a private capability rather than a linear hypothesis.
The monadic presentation itself is standard \cite{wadler1992monads}.

\paragraph{Reference counting and in-place update.}
Lean's runtime is the reference-counting discipline of Ullrich and de Moura
\cite{ullrich2019counting}, with the reuse pass of Perceus
\cite{reinking2021perceus}, and Lean's own machine-checked model of the count
arithmetic is the source of the counter vocabulary in \cref{app:listings}. Koka's
FP${}^2$ \cite{lorenzen2023fp2} is the closest prior art on the other axis: it
statically checks that a function runs \emph{fully} in place. The comparison is
instructive. FP${}^2$ checks the property the compiler needs, by counting, for
arbitrary recursive datatypes, and its check does not model sharing among
readers. Ours models sharing, by construction, and gets one reference at the
write by sealing the state away from clients; it is restricted to a layout
universe of shapes the library can describe. The two are complementary: a
counting discipline answers ``does this function reuse everything'', and a
permission discipline answers ``is this object unshared at this write''.

\paragraph{Uniqueness and linearity in the host language.}
Clean's uniqueness types \cite{barensen1993uniqueness}, Linear Haskell
\cite{bernardy2018linear} and Rust's ownership discipline
\cite{jung2018rustbelt,matsakis2014rust} are the source-language answers to the
same problem: make the count-one condition a typing condition. Lean has no such
answer, which is what makes the interface-based discipline here necessary rather
than merely optional. A reader who has used Rust's borrows will recognise the
shape of what is proved --- a write needs exclusive access, reads may share ---
and will also recognise the difference: Rust's discipline is enforced by the
type checker, ours by a ledger and a sealed interface, and ours therefore has to
say where its enforcement stops.

\paragraph{Unobservable state.}
The pattern of a locally effectful computation observed only through a pure
result is Launchbury and Peyton Jones's \cite{launchbury1995state}: the state
monad's thread is threaded and then hidden, so the effect cannot be observed. Our
situation is different in the mechanism and the same in the spirit. There, the
effect is invisible because the state is not shared; here, the effect is
invisible because the compiler chooses to perform it exactly when no other
reference exists, and the discipline's job is to make that choice reliable.
Reynolds's syntactic control of interference \cite{reynolds1978sci} is the
ancestor of the idea that a program's names, rather than its values, are what
should be controlled; a token that is a name for a ledger entry and not a
permission is a late echo of it.

\paragraph{Levitation and derivatives.}
The description universe is levitation \cite{chapman2010levitation}: describe,
interpret, and prove once. The one-hole context of a path is the derivative of a
data type \cite{mcbride2001derivative}, with the zipper of Huet
\cite{huet1997zipper} as the intuition. Our universe is non-recursive and much
smaller than levitation's; what is carried over is the organisation --- the
description is data the library computes with, and the library, not the author,
writes the operations.

\paragraph{Ghost state and erasure.}
The ledger is ghost state, in the same sense as Dafny's ghost declarations
\cite{leino2010dafny} and F$^{\star}$'s ghost effect \cite{swamy2016fstar}, and the
discipline depends on the compiler erasing arguments by type
\cite{gilbert2019proofirrelevance}. The measurements of what remains at run time
--- that a field-less structure compiles to a constant scalar, that arity
reduction removes literal arguments, that a symbolic argument survives as an
unused borrowed parameter --- are properties of this toolchain's erasure, and
they are the reason the discipline is free rather than merely cheap.

\paragraph{Trusted compilation.}
The last links of \cref{sec:evidence}'s chain are the compiler and the runtime,
which is the position any verified-by-measurement argument must take. CompCert
\cite{leroy2009compcert} is the standard reference for what taking that chain
seriously looks like; the development here takes the weaker option of naming the
assumptions and measuring what it can.

\section{Conclusion}\label{sec:conclusion}

A pure language can update in place because the update cannot be observed. The
condition is that the updated object has one reference, and that condition is a
fact about the compiled code rather than about the meaning of the program, so
getting it where one wants it is a discipline rather than an inherent property
of the source language.

This paper has followed that discipline from its premise to its generalisation.
The premise was made precise in a small counted calculus, where the store rule
at count one is proved equivalent to the copy and the store rule at higher counts
is proved not to be. The discipline was given as fractional permissions in an
erased ghost ledger, behind a program type that no client can build and that one
entry point can run, which makes a write's authority a proof obligation and
costs nothing at run time. It was then lifted three times: to any object, by
letting the object carry its own check; to any state type, by turning the logical
part of the object's rules into proof fields and stating the model's theorems
once; and to any shape, by describing shapes in a universe and transporting the
universe's one instance along a bijection to the author's type. Finally, the
route by which an author could still hand back the state was closed, by a class
in which the library owns the reading, the writing and the encoding, and by a
path calculus in which a whole-array read cannot even be written down.

Two things are, we think, the parts worth keeping. The first is the shape of the
result. Nothing here is a new type system and nothing here changes Lean. The
discipline's enforcement is spread over four mechanisms, each doing what it is
good at: proofs for the accounting, private constructors and a capability for
sealing, an equality between a type and a described layout for the meaning of a
read, and a check that is the identity in the logic for making a copy loud at run
time. The second is the honest accounting. The paper's chain of evidence has
nine links, of which five are proved or structural, two are measured, and two
are assumed; knowing which are which is most of the work of trusting any of it.

The obvious extensions are the ones the design already points at: a recursive
description universe, so that lists and trees are describable and not only
arrays and products; layouts that \emph{name} a user's structure, so that the
equality the sealed class needs can be given for declared types rather than only
for types the layout interpretation denotes definitionally; per-path permissions,
so that two writers can hold disjoint parts of one object; and a mutable cell
behind the sealed interface, which is what a real state variable would want and
which the current development approximates with a logical heap. A concurrent
backend is the largest of them, and \cref{rem:parallel} is the note on why it is
a runtime problem as much as a permission problem.

\appendix
\section{Quoted declarations}\label{app:listings}

The appendix reproduces the definitions and the law statements the paper's
claims rest on. A frame is a line range of the library with its comments
removed --- a frame is code, not documentation --- so what is shown is the
source's own lines rather than a transcription of them, and the range named
beside it is the library's, comments included. The frames are generated by
\code{gen-frames.py} and checked by \code{check-listings.py}, which recomputes
each one from the library and fails when the two differ. The highlighting is
\code{minted}'s, reading each frame through Pygments' Lean lexer.

The description, observation and layout universes are quoted where they are
defined, in \cref{sec:levitation} and \cref{sec:sealed}; what follows is the
rest of the types and the operations, 32 panels, and then 110 law statements,
grouped as the prose groups them. Each frame is one unbreakable box, so a frame
that does not fit the space left on a page moves to the next page instead of
being split across two.

Two further conventions. A frame is a whole source line range, so a declaration
short enough to fit on one line shows its proof on that line: the frame is the
line, not a truncation of it. And a law frame quotes the \emph{statement} of a
theorem, ending where the statement ends, so the proof that follows it is
usually not shown.

The library publishes 184 theorems under \code{Scidonia/}, of which the second
half of this appendix reproduces 110, being the ones this paper states or uses;
seventeen of the 184 are vendored from Lean's own model of the count, and the
seven frames of the reference-count section are this repository's own facts.
Two checks run against the frames. Every panel carries a \code{\%!check:}
directive naming a fragment the quoted range must contain, and every law frame
names its declaration instead; \code{check-listings.py} fails when a range has
drifted either way. The ranges are produced by \code{gen-laws.py}. This checks
that the paper quotes what it claims to quote, and not that the Lean compiles:
the library is built separately.

\newcommand{\listingcaption}[3]{\par\medskip\noindent\textbf{#1}\hfill
  {\footnotesize\itshape\path{#2}, lines #3}\par\nobreak}
\newcommand{\lawframe}[3]{%
  \par\smallskip\noindent\texttt{#3}\hfill
  {\footnotesize\itshape\path{#1}:#2}\par\nobreak
  \codeframe[linerange=#2]{../#1}}

\subsection{The types and the operations}

%!check: def held (l : Ledger) : Rat := 1 - (l.map Prod.snd).sum
\listingcaption{The ledger and the held fraction}{Scidonia/FracPerm/Ledger.lean}{16--22}
\codeframe[linerange=16-22]{../Scidonia/FracPerm/Ledger.lean}

%!check: check_eq : ∀ site o s, check site o s = s
\listingcaption{The object interface}{Scidonia/FracPerm/Hoare.lean}{107--132}
\codeframe[linerange=107-132]{../Scidonia/FracPerm/Hoare.lean}

%!check: structure ReadTok (n : Nat) (q : Rat) : Type where
\listingcaption{The read token}{Scidonia/FracPerm/Hoare.lean}{145--152}
\codeframe[linerange=145-152]{../Scidonia/FracPerm/Hoare.lean}

%!check: private run : Cap → (s : σ) → (∃ l, P l s) →
\listingcaption{The program type}{Scidonia/FracPerm/Hoare.lean}{178--184}
\codeframe[linerange=178-184]{../Scidonia/FracPerm/Hoare.lean}

%!check: def write [FracObject σ ReadOp WriteOp] (o : WriteOp) :
\listingcaption{The primitives and their specifications}{Scidonia/FracPerm/Hoare.lean}{229--266}
\codeframe[linerange=229-266]{../Scidonia/FracPerm/Hoare.lean}

%!check: inductive WriteOp where
\listingcaption{The second object: a count and a log}{Scidonia/FracPerm/Counter.lean}{22--49}
\codeframe[linerange=22-49]{../Scidonia/FracPerm/Counter.lean}

%!check: inductive Value where
\listingcaption{The value language, in which a reference can be written down}{Scidonia/FracPerm/StateVar/Value.lean}{46--67}
\codeframe[linerange=46-67]{../Scidonia/FracPerm/StateVar/Value.lean}

%!check: def NoRef (ℓ : Loc) : Value → Prop
\listingcaption{Detachment}{Scidonia/FracPerm/StateVar/Value.lean}{68--82}
\codeframe[linerange=68-82]{../Scidonia/FracPerm/StateVar/Value.lean}

%!check: observe_detached : ∀ (ℓ : Loc) (o : ReadOp) (s : σ),
\listingcaption{The lawful state variable}{Scidonia/FracPerm/StateVar/Class.lean}{59--91}
\codeframe[linerange=59-91]{../Scidonia/FracPerm/StateVar/Class.lean}

%!check: abbrev Heap (σ : Type) := Loc → Option σ
\listingcaption{The single-root heap}{Scidonia/FracPerm/StateVar/Heap.lean}{44--47}
\codeframe[linerange=44-47]{../Scidonia/FracPerm/StateVar/Heap.lean}

%!check: def Heap.update (h : Heap σ) (ℓ : Loc) (v : Option σ) : Heap σ :=
\listingcaption{The heap's own point update and its two laws}{Scidonia/FracPerm/StateVar/Heap.lean}{49--68}
\codeframe[linerange=49-68]{../Scidonia/FracPerm/StateVar/Heap.lean}

%!check: def ReadAt {σ Init ReadOp WriteOp : Type} (inst : LawfulStateVar σ Init ReadOp WriteOp)
\listingcaption{A read step against the heap}{Scidonia/FracPerm/StateVar/Heap.lean}{70--77}
\codeframe[linerange=70-77]{../Scidonia/FracPerm/StateVar/Heap.lean}

%!check: def WriteAt {σ Init ReadOp WriteOp : Type} (inst : LawfulStateVar σ Init ReadOp WriteOp)
\listingcaption{A write step against the heap}{Scidonia/FracPerm/StateVar/Heap.lean}{79--89}
\codeframe[linerange=79-89]{../Scidonia/FracPerm/StateVar/Heap.lean}

%!check: inductive RTree (ReadOp : Type) (Result : ReadOp → Type) : Type → Type 1 where
\listingcaption{The generic recipes}{Scidonia/FracPerm/StateVar/Surface.lean}{38--59}
\codeframe[linerange=38-59]{../Scidonia/FracPerm/StateVar/Surface.lean}

%!check: def arrayObserve (o : ArrayReadOp) (a : Array UInt64) : ArrayResult o :=
\listingcaption{The array witness}{Scidonia/FracPerm/StateVar/Array.lean}{57--76}
\codeframe[linerange=57-76]{../Scidonia/FracPerm/StateVar/Array.lean}

%!check: structure Described (d : Desc) where
\listingcaption{The description-indexed wrapper}{Scidonia/FracPerm/StateVar/Universe.lean}{111--117}
\codeframe[linerange=111-117]{../Scidonia/FracPerm/StateVar/Universe.lean}

%!check: LawfulStateVar τ τ Unit τ where
\listingcaption{Transport along an isomorphism}{Scidonia/FracPerm/StateVar/Universe.lean}{210--240}
\codeframe[linerange=210-240]{../Scidonia/FracPerm/StateVar/Universe.lean}

%!check: def encode : Slot → El slotDesc
\listingcaption{The slot, and its isomorphism to its description}{Scidonia/FracPerm/StateVar/Slot.lean}{46--63}
\codeframe[linerange=46-63]{../Scidonia/FracPerm/StateVar/Slot.lean}


%!check: rep : σ = State layout
\listingcaption{The sealed class}{Scidonia/FracPerm/StateVar/Sealed.lean}{52--60}
\codeframe[linerange=52-60]{../Scidonia/FracPerm/StateVar/Sealed.lean}

%!check: fun s => Layout5.read? p (toState s)
\listingcaption{The library's read of a sealed state}{Scidonia/FracPerm/StateVar/Sealed.lean}{66--74}
\codeframe[linerange=66-74]{../Scidonia/FracPerm/StateVar/Sealed.lean}

%!check: (ofState r.1, r.2)
\listingcaption{The library's write of a sealed state}{Scidonia/FracPerm/StateVar/Sealed.lean}{83--90}
\codeframe[linerange=83-90]{../Scidonia/FracPerm/StateVar/Sealed.lean}

%!check: def Disjoint {d : Layout} (p q : Path d) : Prop :=
\listingcaption{Disjointness of paths}{Scidonia/FracPerm/StateVar/Layout5.lean}{203--208}
\codeframe[linerange=203-208]{../Scidonia/FracPerm/StateVar/Layout5.lean}

%!check: private def takeOut (s : Array α) (i : Nat) (h : i < s.size) (d : α) : Array α × α
\listingcaption{The take-out}{Scidonia/FracPerm/StateVar/Layout5.lean}{298--308}
\codeframe[linerange=298-308]{../Scidonia/FracPerm/StateVar/Layout5.lean}

%!check: def encodeObs : (o : ObsDesc) → ObsEl o → Value
\listingcaption{The library's encoding of an observation}{Scidonia/FracPerm/StateVar/Layout5.lean}{457--471}
\codeframe[linerange=457-471]{../Scidonia/FracPerm/StateVar/Layout5.lean}

%!check: private inductive Context : (d : Layout) → Path d → Type where
\listingcaption{The one-hole context}{Scidonia/FracPerm/StateVar/Layout5.lean}{539--555}
\codeframe[linerange=539-555]{../Scidonia/FracPerm/StateVar/Layout5.lean}

%!check: def read? {d : Layout} (p : ObsPath d) (s : State d) : Option Value
\listingcaption{The library's read}{Scidonia/FracPerm/StateVar/Layout5.lean}{679--684}
\codeframe[linerange=679-684]{../Scidonia/FracPerm/StateVar/Layout5.lean}

%!check: private def setAt : {d : Layout} → (p : Path d) → State (At p) → State d → State d
\listingcaption{The fused update}{Scidonia/FracPerm/StateVar/Layout5.lean}{710--729}
\codeframe[linerange=710-729]{../Scidonia/FracPerm/StateVar/Layout5.lean}

%!check: (setAt p.path (cast (congrArg State p.leaf.symm) x) s, true)
\listingcaption{The public write}{Scidonia/FracPerm/StateVar/Layout5.lean}{772--784}
\codeframe[linerange=772-784]{../Scidonia/FracPerm/StateVar/Layout5.lean}

%!check: def IsExclusive (rc : Int32) : Prop := isSt rc ∧ rc = 1
\listingcaption{What a write tests}{Scidonia/FracPerm/StateVar/RefCount.lean}{713--728}
\codeframe[linerange=713-728]{../Scidonia/FracPerm/StateVar/RefCount.lean}

%!check: inductive RTree : Type → Type 1 where
\listingcaption{The structured interface's command trees}{Scidonia/FracPerm/Structured/Surface.lean}{127--149}
\codeframe[linerange=127-149]{../Scidonia/FracPerm/Structured/Surface.lean}

%!check: structure Reader (α : Type) where
\listingcaption{The reader and the program, each behind a private field}{Scidonia/FracPerm/Structured/Surface.lean}{166--179}
\codeframe[linerange=166-179]{../Scidonia/FracPerm/Structured/Surface.lean}

%!check: structure Init where
\listingcaption{What a fresh session starts from}{Scidonia/FracPerm/Structured/Surface.lean}{263--270}
\codeframe[linerange=263-270]{../Scidonia/FracPerm/Structured/Surface.lean}

% Generated from the library by gen-laws.py; do not edit by hand.
% Each \lawframe{module}{lines}{name} quotes the statement of `name` from
% that module.  check-listings.py verifies that every range still holds the
% declaration it names, so the appendix cannot drift from the library.

\subsection{The ledger and the arithmetic of shares}
\lawframe{Scidonia/FracPerm/Ledger.lean}{25-25}{held\_nil}
\lawframe{Scidonia/FracPerm/Ledger.lean}{28-28}{held\_cons}
\lawframe{Scidonia/FracPerm/Ledger.lean}{34-35}{held\_erase\_of\_mem}

\subsection{Detachment in the value language}
\lawframe{Scidonia/FracPerm/StateVar/Value.lean}{84-84}{noRef\_loc}
\lawframe{Scidonia/FracPerm/StateVar/Value.lean}{96-97}{noRef\_pair}
\lawframe{Scidonia/FracPerm/StateVar/Value.lean}{109-109}{noRef\_loc\_of\_ne}
\lawframe{Scidonia/FracPerm/StateVar/Value.lean}{113-113}{not\_noRef\_loc\_self}

\subsection{The generic model}
\lawframe{Scidonia/FracPerm/StateVar/Model.lean}{369-370}{primitive\_detached}
\lawframe{Scidonia/FracPerm/StateVar/Model.lean}{423-424}{replay\_trace}
\lawframe{Scidonia/FracPerm/StateVar/Model.lean}{522-523}{eval\_of\_run}
\lawframe{Scidonia/FracPerm/StateVar/Model.lean}{529-530}{complete\_run\_value}
\lawframe{Scidonia/FracPerm/StateVar/Model.lean}{535-535}{wf\_leafSum}
\lawframe{Scidonia/FracPerm/StateVar/Model.lean}{550-551}{wf\_sharesPositive}
\lawframe{Scidonia/FracPerm/StateVar/Model.lean}{597-600}{reading\_step\_is\_read\_or\_leave}
\lawframe{Scidonia/FracPerm/StateVar/Model.lean}{607-608}{reading\_step\_store}
\lawframe{Scidonia/FracPerm/StateVar/Model.lean}{582-587}{wait\_step}
\lawframe{Scidonia/FracPerm/StateVar/Model.lean}{642-643}{writer\_authority\_add}
\lawframe{Scidonia/FracPerm/StateVar/Model.lean}{656-657}{pstep\_wf}
\lawframe{Scidonia/FracPerm/StateVar/Model.lean}{667-668}{safe\_result}

\subsection{The single-root heap}
\lawframe{Scidonia/FracPerm/StateVar/Heap.lean}{102-103}{read\_at\_detached}
\lawframe{Scidonia/FracPerm/StateVar/Heap.lean}{118-119}{write\_at\_valid}
\lawframe{Scidonia/FracPerm/StateVar/Heap.lean}{126-127}{write\_at\_frame}
\lawframe{Scidonia/FracPerm/StateVar/Heap.lean}{135-140}{snapshot\_after\_write}
\lawframe{Scidonia/FracPerm/StateVar/Heap.lean}{157-162}{pstep\_write\_corresponds}
\lawframe{Scidonia/FracPerm/StateVar/Heap.lean}{168-171}{readAt\_step\_observe}
\lawframe{Scidonia/FracPerm/StateVar/Heap.lean}{214-214}{cannot\_add\_reader}

\subsection{The array witness and the bridge}
\lawframe{Scidonia/FracPerm/StateVar/Array.lean}{80-81}{arrayObserve\_detached}
\lawframe{Scidonia/FracPerm/StateVar/Array.lean}{91-92}{encodeArrayRead\_injective}
\lawframe{Scidonia/FracPerm/StateVar/Array.lean}{113-114}{arrayReplace\_size}
\lawframe{Scidonia/FracPerm/StateVar/Array.lean}{120-121}{arrayReplace\_bounds}
\lawframe{Scidonia/FracPerm/StateVar/Array.lean}{165-166}{bridge\_val}
\lawframe{Scidonia/FracPerm/StateVar/Array.lean}{194-195}{bridge\_pval}
\lawframe{Scidonia/FracPerm/StateVar/Array.lean}{216-220}{runFresh\_eq\_pval}

\subsection{The description universe and transport}
\lawframe{Scidonia/FracPerm/StateVar/Universe.lean}{86-87}{encodeEl\_injective}
\lawframe{Scidonia/FracPerm/StateVar/Universe.lean}{103-103}{encodeEl\_noRef}
\lawframe{Scidonia/FracPerm/StateVar/Universe.lean}{153-154}{universe\_read\_after\_replace}
\lawframe{Scidonia/FracPerm/StateVar/Universe.lean}{159-160}{universe\_read\_detached}
\lawframe{Scidonia/FracPerm/StateVar/Universe.lean}{178-181}{universe\_write\_at\_valid}
\lawframe{Scidonia/FracPerm/StateVar/Universe.lean}{185-187}{universe\_write\_at\_frame}
\lawframe{Scidonia/FracPerm/StateVar/Universe.lean}{193-199}{universe\_snapshot\_after\_write}
\lawframe{Scidonia/FracPerm/StateVar/Universe.lean}{273-277}{transport\_read\_after\_replace}

\subsection{The slot}
\lawframe{Scidonia/FracPerm/StateVar/Slot.lean}{67-67}{decode\_encode}
\lawframe{Scidonia/FracPerm/StateVar/Slot.lean}{73-73}{encode\_decode}
\lawframe{Scidonia/FracPerm/StateVar/Slot.lean}{82-82}{lifted\_iff}
\lawframe{Scidonia/FracPerm/StateVar/Slot.lean}{131-132}{slot\_read\_after\_replace}
\lawframe{Scidonia/FracPerm/StateVar/Slot.lean}{152-156}{slot\_snapshot\_full7\_full9}
\lawframe{Scidonia/FracPerm/StateVar/Slot.lean}{169-174}{slot\_shape\_change}
\lawframe{Scidonia/FracPerm/StateVar/Slot.lean}{192-192}{slot\_full\_round\_trip}

\subsection{Layouts, paths and the take-out}
\lawframe{Scidonia/FracPerm/StateVar/Layout5.lean}{190-191}{append\_left}
\lawframe{Scidonia/FracPerm/StateVar/Layout5.lean}{194-195}{append\_right}
\lawframe{Scidonia/FracPerm/StateVar/Layout5.lean}{198-199}{append\_index}
\lawframe{Scidonia/FracPerm/StateVar/Layout5.lean}{210-211}{not\_disjoint\_here\_left}
\lawframe{Scidonia/FracPerm/StateVar/Layout5.lean}{221-222}{disjoint\_left\_iff}
\lawframe{Scidonia/FracPerm/StateVar/Layout5.lean}{270-271}{disjoint\_index\_ne}
\lawframe{Scidonia/FracPerm/StateVar/Layout5.lean}{474-475}{encodeObs\_injective}
\lawframe{Scidonia/FracPerm/StateVar/Layout5.lean}{511-512}{encodeObs\_noRef}
\lawframe{Scidonia/FracPerm/StateVar/Layout5.lean}{1027-1028}{setLeaf\_failed}
\lawframe{Scidonia/FracPerm/StateVar/Layout5.lean}{1055-1056}{read\_detached}
\lawframe{Scidonia/FracPerm/StateVar/Layout5.lean}{1243-1246}{snapshot\_after\_setLeaf}
\lawframe{Scidonia/FracPerm/StateVar/Layout5.lean}{1265-1267}{read\_frame}
\lawframe{Scidonia/FracPerm/StateVar/Layout5.lean}{1283-1284}{array\_root\_not\_observable}
\lawframe{Scidonia/FracPerm/StateVar/Layout5.lean}{1361-1364}{describedLayout\_read\_after\_replace}
\lawframe{Scidonia/FracPerm/StateVar/Layout5.lean}{1379-1382}{describedLayout\_primitive\_detached}
\lawframe{Scidonia/FracPerm/StateVar/Layout5.lean}{1402-1405}{describedLayout\_write\_at\_frame}

\subsection{The sealed class}
\lawframe{Scidonia/FracPerm/StateVar/Sealed.lean}{125-127}{sealedLawful\_Result}
\lawframe{Scidonia/FracPerm/StateVar/Sealed.lean}{132-135}{sealedLawful\_encodeRead}
\lawframe{Scidonia/FracPerm/StateVar/Sealed.lean}{152-153}{toState\_ofState}
\lawframe{Scidonia/FracPerm/StateVar/Sealed.lean}{157-158}{ofState\_toState}
\lawframe{Scidonia/FracPerm/StateVar/Sealed.lean}{191-193}{sealedRead\_detached}
\lawframe{Scidonia/FracPerm/StateVar/Sealed.lean}{199-202}{sealedRead\_after\_setLeaf}
\lawframe{Scidonia/FracPerm/StateVar/Sealed.lean}{213-217}{sealedSnapshot}
\lawframe{Scidonia/FracPerm/StateVar/Sealed.lean}{234-237}{sealedRead\_frame}
\lawframe{Scidonia/FracPerm/StateVar/Sealed.lean}{251-253}{sealedSetLeaf\_failed}

\subsection{The reference count}
\lawframe{Scidonia/FracPerm/StateVar/RefCount.lean}{740-741}{not\_exclusive\_incRefN\_one}
\lawframe{Scidonia/FracPerm/StateVar/RefCount.lean}{765-765}{decRef\_two}
\lawframe{Scidonia/FracPerm/StateVar/RefCount.lean}{771-771}{not\_exclusive\_of\_isMt}
\lawframe{Scidonia/FracPerm/StateVar/RefCount.lean}{779-779}{markMtRc\_not\_exclusive}
\lawframe{Scidonia/FracPerm/StateVar/RefCount.lean}{787-788}{not\_exclusive\_of\_neverFreed}
\lawframe{Scidonia/FracPerm/StateVar/RefCount.lean}{804-805}{run\_never\_exclusive}
\lawframe{Scidonia/FracPerm/StateVar/RefCount.lean}{819-823}{takeOut\_refCount\_eq\_one}

\subsection{The structured model}
\lawframe{Scidonia/FracPerm/Structured/Model.lean}{292-293}{step\_get}
\lawframe{Scidonia/FracPerm/Structured/Model.lean}{297-298}{step\_size}
\lawframe{Scidonia/FracPerm/Structured/Model.lean}{203-204}{replay\_trace}
\lawframe{Scidonia/FracPerm/Structured/Model.lean}{365-366}{eval\_of\_run}
\lawframe{Scidonia/FracPerm/Structured/Model.lean}{372-373}{complete\_run\_value}
\lawframe{Scidonia/FracPerm/Structured/Model.lean}{415-415}{wf\_leafSum}
\lawframe{Scidonia/FracPerm/Structured/Model.lean}{430-431}{wf\_sharesPositive}
\lawframe{Scidonia/FracPerm/Structured/Model.lean}{549-552}{reading\_step\_is\_read\_or\_leave}
\lawframe{Scidonia/FracPerm/Structured/Model.lean}{559-560}{reading\_step\_store}
\lawframe{Scidonia/FracPerm/Structured/Model.lean}{463-468}{wait\_step}
\lawframe{Scidonia/FracPerm/Structured/Model.lean}{565-567}{set\_step}
\lawframe{Scidonia/FracPerm/Structured/Model.lean}{572-573}{set\_step\_size}
\lawframe{Scidonia/FracPerm/Structured/Model.lean}{586-587}{writer\_authority\_add}

\subsection{What the permission discipline refuses}
\lawframe{ScidoniaTest/FracPerm.lean}{96-96}{n1\_write\_while\_read\_outstanding}
\lawframe{ScidoniaTest/FracPerm.lean}{110-110}{n2\_release\_twice}
\lawframe{ScidoniaTest/FracPerm.lean}{124-124}{n3\_read\_with\_returned\_token}
\lawframe{ScidoniaTest/FracPerm.lean}{136-136}{n4\_overcommit}
\lawframe{ScidoniaTest/FracPerm.lean}{149-149}{n4b\_reused\_ghost\_id}
\lawframe{ScidoniaTest/FracPerm.lean}{160-161}{n5\_unreturned\_read}
\lawframe{ScidoniaTest/FracPerm.lean}{313-313}{w3\_escaped\_token}

\subsection{Countermodels: the laws and the count are not vacuous}
\lawframe{ScidoniaTest/StateVar.lean}{281-281}{bad\_read\_rejected}
\lawframe{ScidoniaTest/StateVar.lean}{287-287}{bad\_write\_rejected}
\lawframe{ScidoniaTest/StateVar.lean}{293-293}{bad\_encode\_rejected}
\lawframe{ScidoniaTest/Model5.lean}{335-336}{leakyArray\_read\_returns\_state}
\lawframe{ScidoniaTest/Model5.lean}{353-353}{leaky\_root\_not\_detached}
\lawframe{ScidoniaTest/Count.lean}{44-44}{one\_is\_exclusive}
\lawframe{ScidoniaTest/Count.lean}{56-57}{stuckSt\_is\_unshared\_not\_exclusive}
\lawframe{ScidoniaTest/Count.lean}{72-73}{shared\_count\_is\_not\_exclusive}
\lawframe{ScidoniaTest/Count.lean}{125-125}{not\_exclusive\_after\_handOff}
\lawframe{ScidoniaTest/Count.lean}{133-135}{takeOut\_counts\_at\_one}


\begin{thebibliography}{99}\small

\bibitem{barensen1993uniqueness}
E.~Barendsen and S.~Smetsers.
\newblock Conventional and uniqueness typing in graph rewrite systems.
\newblock In \emph{FSTTCS}, 1993.

\bibitem{bernardy2018linear}
J.-P.~Bernardy, M.~Boespflug, R.~R.~Newton, S.~Peyton Jones, and A.~Spiwack.
\newblock Linear Haskell: practical linearity in a higher-order polymorphic
language.
\newblock \emph{POPL}, 2018.

\bibitem{boyland2003checking}
J.~Boyland.
\newblock Checking interference with fractional permissions.
\newblock In \emph{Static Analysis (SAS)}, LNCS 2694, pages 55--72, 2003.
\newblock \doi{10.1007/3-540-44898-5\_4}.

\bibitem{chapman2010levitation}
J.~Chapman, P.-E.~Dagand, C.~McBride, and P.~Morris.
\newblock The gentle art of levitation.
\newblock In \emph{ICFP}, pages 3--14, 2010.
\newblock \doi{10.1145/1863543.1863547}.

\bibitem{fromherz2021steel}
A.~Fromherz, A.~Rastogi, N.~Swamy, S.~Gibson, G.~Mart\'inez, D.~Merigoux, and
T.~Ramananandro.
\newblock Steel: proof-oriented programming in a dependently typed concurrent
separation logic.
\newblock In \emph{ICFP}, 2021.
\newblock \doi{10.1145/3473590}.

\bibitem{gilbert2019proofirrelevance}
G.~Gilbert, J.~Cockx, M.~Sozeau, and N.~Tabareau.
\newblock Definitional proof-irrelevance without K.
\newblock \emph{POPL}, 2019.

\bibitem{huet1997zipper}
G.~Huet.
\newblock The zipper.
\newblock \emph{Journal of Functional Programming}, 7(5):549--554, 1997.

\bibitem{jacobs2011verifast}
B.~Jacobs, J.~Smans, P.~Philippaerts, F.~Vogels, W.~Penninckx, and
F.~Piessens.
\newblock VeriFast: a powerful, sound, predictable, fast verifier for C and
Java.
\newblock In \emph{NASA Formal Methods}, 2011.

\bibitem{jung2015iris}
R.~Jung, D.~Swasey, F.~Sieczkowski, K.~Svendsen, A.~Turon, L.~Birkedal, and
D.~Dreyer.
\newblock Iris: monoids and invariants as an orthogonal basis for concurrent
reasoning.
\newblock \emph{POPL}, 2015.

\bibitem{jung2018rustbelt}
R.~Jung, J.-H.~Jourdan, R.~Krebbers, and D.~Dreyer.
\newblock RustBelt: securing the foundations of the Rust programming language.
\newblock \emph{POPL}, 2018.

\bibitem{krebbers2017essence}
R.~Krebbers, A.~Timany, and L.~Birkedal.
\newblock The essence of higher-order concurrent separation logic.
\newblock In \emph{ESOP}, 2017.

\bibitem{lattuada2023verus}
A.~Lattuada, T.~Hance, C.~Cho, M.~Brun, I.~Subasinghe, Y.~Zhou, J.~Howell,
B.~Parno, and C.~Hawblitzel.
\newblock Verus: verifying Rust programs using linear ghost types.
\newblock \emph{OOPSLA}, 2023.
\newblock \doi{10.1145/3586037}.

\bibitem{launchbury1995state}
J.~Launchbury and S.~Peyton Jones.
\newblock State in Haskell.
\newblock \emph{Lisp and Symbolic Computation}, 8(4):293--341, 1995.

\bibitem{lean4rcmodel}
Lean FRO.
\newblock \code{tests/elab/rc\_model.lean}, Lean~4, tag \code{v4.35.0-rc3}.
\newblock A machine-checked model of \code{lean\_inc\_ref\_n} and
\code{lean\_dec\_ref}, vendored into \code{Scidonia/FracPerm/StateVar/RefCount.lean}.

\bibitem{leino2010dafny}
K.~R.~M.~Leino.
\newblock Dafny: an automatic program verifier for functional correctness.
\newblock In \emph{LPAR}, 2010.

\bibitem{leroy2009compcert}
X.~Leroy.
\newblock Formal verification of a realistic compiler.
\newblock \emph{Communications of the ACM}, 52(7):107--115, 2009.

\bibitem{lorenzen2023fp2}
A.~Lorenzen, D.~Leijen, and W.~Swierstra.
\newblock FP${}^2$: fully in-place functional programming.
\newblock \emph{PACMPL}, 7(ICFP), 2023.
\newblock \doi{10.1145/3607840}.

\bibitem{matsakis2014rust}
N.~D.~Matsakis and F.~S.~Klock.
\newblock The Rust language.
\newblock \emph{ACM SIGAda Ada Letters}, 34(3), 2014.

\bibitem{mcbride2001derivative}
C.~McBride.
\newblock The derivative of a regular type is its type of one-hole contexts.
\newblock Extended abstract, 2001.

\bibitem{muller2016viper}
P.~M\"uller, M.~Schwerhoff, and A.~J.~Summers.
\newblock Viper: a verification infrastructure for permission-based reasoning.
\newblock In \emph{VMCAI}, 2016.

\bibitem{nanevski2008htt}
A.~Nanevski, G.~Morrisett, and L.~Birkedal.
\newblock Hoare type theory, polymorphism and separation.
\newblock \emph{Journal of Functional Programming}, 18(5--6), 2008.

\bibitem{nanevski2008ynot}
A.~Nanevski, G.~Morrisett, A.~Shinnar, P.~Govereau, and L.~Birkedal.
\newblock Ynot: dependent types for imperative programs.
\newblock In \emph{ICFP}, 2008.

\bibitem{ohearn2001local}
P.~W.~O'Hearn, J.~C.~Reynolds, and H.~Yang.
\newblock Local reasoning about programs that alter data structures.
\newblock In \emph{Computer Science Logic (CSL)}, 2001.

\bibitem{reinking2021perceus}
A.~Reinking, N.~Xie, L.~de~Moura, and D.~Leijen.
\newblock Perceus: garbage free reference counting with reuse.
\newblock In \emph{PLDI}, 2021.
\newblock \doi{10.1145/3453483.3454032}.

\bibitem{reynolds1978sci}
J.~C.~Reynolds.
\newblock Syntactic control of interference.
\newblock In \emph{POPL}, 1978.

\bibitem{reynolds2002separation}
J.~C.~Reynolds.
\newblock Separation logic: a logic for shared mutable data structures.
\newblock In \emph{LICS}, 2002.

\bibitem{swamy2016fstar}
N.~Swamy, C.~Hri\c{t}cu, C.~Keller, A.~Rastogi, A.~Delignat-Lavaud, S.~Forest,
K.~Bhargavan, C.~Fournet, P.-Y.~Strub, M.~Kohlweiss, J.-K.~Zinzindohou\'e, and
S.~Zanella-B\'eguelin.
\newblock Dependent types and multi-monadic effects in F$^{\star}$.
\newblock In \emph{POPL}, 2016.

\bibitem{swierstra2009hoare}
W.~Swierstra.
\newblock A Hoare logic for the state monad.
\newblock In \emph{Theorem Proving in Higher Order Logics (TPHOLs)}, 2009.

\bibitem{ullrich2019counting}
S.~Ullrich and L.~de~Moura.
\newblock Counting immutable beans: reference counting optimized for purely
functional programming.
\newblock In \emph{Implementation and Application of Functional Languages
(IFL)}, 2019.
\newblock \doi{10.1145/3412932.3412935}.

\bibitem{wadler1992monads}
P.~Wadler.
\newblock The essence of functional programming.
\newblock In \emph{POPL}, 1992.

\end{thebibliography}

\end{document}
